\documentclass[11pt,a4paper]{article}
\usepackage[T1]{fontenc}
\usepackage[utf8]{inputenc}
\usepackage{lmodern,microtype}
\usepackage[margin=.85in]{geometry}
\usepackage{amsmath,amssymb,amsthm,mathtools,aliascnt}
\usepackage{booktabs,tabularx,array,enumitem,needspace,float}
\usepackage[section]{placeins}
\usepackage{xcolor,tikz,mdframed}
\usetikzlibrary{arrows.meta,positioning,calc,decorations.pathreplacing}
\usepackage[hidelinks]{hyperref}
\usepackage[nameinlink,noabbrev]{cleveref}
\usepackage{fancyhdr}
\hypersetup{pdftitle={Sync HotStuff RR, Part II v2: Sequencer Rotation},pdfsubject={Espresso-confirmed rotation with automatic and multisignature Force}}
\setlength{\headheight}{14pt}
\pagestyle{fancy}\fancyhf{}
\lhead{\small Sync HotStuff RR, Part II}\rhead{\small Version 2 / Espresso finality}
\cfoot{\thepage}
\setlength{\parindent}{0pt}\setlength{\parskip}{4pt}
\setlength{\emergencystretch}{2em}
\renewcommand{\topfraction}{.92}
\renewcommand{\textfraction}{.06}
\renewcommand{\floatpagefraction}{.85}
\makeatletter
\setlength{\@fptop}{0pt}\setlength{\@fpbot}{0pt plus 1fil}
\makeatother
\setlist{nosep,leftmargin=*}
\newcolumntype{Y}{>{\raggedright\arraybackslash}X}
\definecolor{rrRed}{HTML}{A51C30}
\definecolor{rrBlue}{HTML}{2563EB}
\definecolor{rrTeal}{HTML}{0F766E}
\newcommand{\RR}[1]{{\color{rrRed}#1}}
\newcommand{\QC}{\mathcal C}
\newcommand{\FC}{\mathcal F}
\newcommand{\Hash}{\mathsf H}
\newcommand{\rank}{\operatorname{rank}}
\newcommand{\TB}{\mathit{TB}}
\newcommand{\tagg}[1]{\ensuremath{\mathsf{#1}}}
\newcommand{\defterm}[1]{\textcolor{rrBlue}{\textit{#1}}}
\newtheorem{definition}{Definition}[section]
\newaliascnt{assumption}{definition}
\newtheorem{assumption}[assumption]{Assumption}
\aliascntresetthe{assumption}
\crefname{assumption}{assumption}{assumptions}
\Crefname{assumption}{Assumption}{Assumptions}
\newcounter{protocol}
% Protocol boxes float to the top of a page, so long boxes leave no gaps.
\newenvironment{protocol}[2]{\begin{figure}[htbp]\refstepcounter{protocol}\label{#1}
\begin{mdframed}[linewidth=.6pt,skipabove=0pt,skipbelow=0pt,
innertopmargin=7pt,innerbottommargin=7pt]\small
\textbf{Protocol \theprotocol. #2}\par\smallskip}{\end{mdframed}\end{figure}}
\crefname{protocol}{Protocol}{Protocols}
\Crefname{protocol}{Protocol}{Protocols}
\title{Sync HotStuff RR, Part II: Sequencer Rotation\\\large Version 2: Espresso-Confirmed Multi-Region Rotation}
\author{}\date{17 October 2026}
\begin{document}
\maketitle
\begin{abstract}
Sync HotStuff RR lets several regions take turns producing one chain.
Each region runs the Part~I engine with a fixed voting snapshot and a
fixed quota of DATA blocks. A terminal regional receipt opens the next
normal session without waiting for global confirmation. Publishers send
complete regional results to Espresso. An observer confirms a result
globally when it accepts that result as a continuous extension in
Espresso's finalized history. External control uses Espresso timestamps.
A block-missing Force follows from insufficient accepted DATA in a fixed
window. A multisignature Force is announced on L1 before its scheduled
Espresso time. Both preserve the entire global prefix and close the old
production branch to further global output. The regional DATA/DRAIN
engine, its timers, and its recovery rules remain unchanged. We specify
the interfaces, event order, and a conditional prefix-safety argument.
The progress conditions are stated separately.
\end{abstract}

\section{Overview}
\label{sec:overview}
Part~I~\cite{part1} gives one fixed committee a synchronous BFT engine.
This document composes regional instances of that engine. A
\defterm{session} $r$ is one production turn of region $i_r$.
Its view is $(r,g,t)$: session, recovery generation, and leader term.
The voting snapshot and DATA quota are fixed throughout the session.

\paragraph{What this version changes.}
\begin{itemize}
\item \textbf{Normal rotation remains fast.} A valid receipt for the
terminal DATA block $T_r$ opens the normal successor. The successor
needs the boundary, state and authority, but does not wait for $T_r$
to become globally final.
\item \textbf{Espresso supplies global order.} A publisher may submit
a complete range of regional results. Every verifier uses the same
Espresso history and control rules to accept a continuous DATA prefix.
There is no per-block committee acknowledgement or global waiting timer.
\item \textbf{There are two Forces.} An automatic Force follows from
a failed DATA-delivery window. A multisignature Force follows from an
authorized L1 record with a future Espresso activation time.
\item \textbf{A Force preserves all global output.} It closes the old
\defterm{branch epoch} $\beta$ and starts a new one from the current
global prefix. Normal successors on the old epoch cannot publish its
unconfirmed suffix after closure.
\item \textbf{Control is interpreted by verifiers.} Espresso stores
ordered bytes. Nodes and certifiers verify L1 and Espresso evidence and
run the same external interpreter. Espresso need not execute a contract
for this chain.
\end{itemize}

The regional engine retains DATA, recovery-only DRAIN, FINAL receipts,
and Completion. A DRAIN has no application payload and never enters the
global DATA prefix. A Force does not cancel a regional timer, change a
lock, or create a quit proof.

\begin{table}[htbp]\centering\small
\caption{Participants and responsibilities.}\label{tab:roles}
\begin{tabularx}{\linewidth}{@{}>{\raggedright\arraybackslash}p{.19\linewidth}Y@{}}\toprule
Role & Responsibility\\\midrule
L1 & Publishes immutable authorities and the ordered control log.\\
Regional replica & Runs the DATA/DRAIN engine of its sessions and
serves complete commitment evidence.\\
Multisignature service & Announces future status changes and Forces
on L1. It cannot choose a retained prefix or declare data global.\\
Publisher & Obtains complete regional bundles, submits them to
Espresso, and retries ranges that have not been accepted.\\
Espresso & Supplies a finalized ordered log, authenticated timestamps,
and proofs of its complete namespace contents.\\
Observer / certifier & Verifies both sources, replays control and
exports, and derives global output and forced boundaries.\\
Backbone & Eventually delivers regional data, receipts and witnesses.
It does not choose a region or grant finality.\\\bottomrule
\end{tabularx}
\end{table}

\paragraph{Production, confirmation, and publication.}
The \defterm{production branch} $R_u$ of observer $u$ may contain
several normal sessions beyond its \defterm{global prefix} $G_u$.
The \defterm{export session} is the session responsible for the next
DATA height after $G_u$. It need not be the session currently producing
new blocks. The accepted DA publication prefix $C_{\rm pub}$ may lag
behind $G_u$ as well. Neither lag permits a Force to remove $G_u$.

A block's regional receipt remains evidence of its regional commitment
even if its unconfirmed branch is later closed. Global finality is
stronger: accepted DATA is never removed by either Force.
The claim concerns input content, order and provenance. Execution
correctness and native state continuity remain obligations of the
application and publication interfaces.

\begin{table}[htbp]\centering\small
\caption{Changes of production view.}\label{tab:views}
\begin{tabularx}{\linewidth}{@{}>{\raggedright\arraybackslash}p{.19\linewidth}>{\raggedright\arraybackslash}p{.18\linewidth}Y@{}}\toprule
Change & New view & Trigger and timer effect\\\midrule
Normal leader rotation & $(r,g,t+1)$ & A terminal QC below $H_r$;
existing commit timers continue.\\
Regional recovery & $(r,g+1,0)$ & A quit proof, the $\Delta_i$ wait
and finalized authority; the quit cancels generation-$g$ timers.\\
Normal regional handoff & $(r+1,0,0)$ & Ready on a REGULAR witness
with the terminal receipt; the old engine stops through Completion.\\
Forced handoff & $(r_c+1,0,0)$ & Ready on a derived Force boundary;
$r_c$ is the closed export-session index. Old engines are unchanged.\\
\bottomrule\end{tabularx}
\end{table}

A Force can reuse an index from a discarded provisional successor.
Its new epoch and session ID keep all votes, locks and evidence separate.
Normal handoffs retain the epoch. A finalized but not yet effective L1
control record does not stop production or global acceptance.

\section{Model and assumptions}
\label{sec:model}
\subsection{L1 control state}
\label{sec:l1}
Nodes authenticate finalized L1 headers, authorities and a complete
ordered control prefix. L1 timestamps label L1 blocks; they are not
regional commit times or Espresso activation times.

\begin{table}[htbp]\centering\small
\caption{L1 authority and control records.}\label{tab:l1}
\begin{tabularx}{\linewidth}{@{}p{.27\linewidth}Y@{}}\toprule
Field & Meaning\\\midrule
$\tagg{AUTH}_{r,0}$ & The fixed regional snapshots, leader schedule,
DATA quota $\TB_r$, and export-window parameters for index $r$.\\
$\tagg{AUTH}_{r,g}$, $g>0$ & The leader schedule for recovery
generation $g$; it does not change the original snapshot or quota.\\
Control sequence $c$ & A unique increasing index in the L1 control
log. It binds the chain, record kind and payload.\\
Activation time $\theta_a$ & An Espresso timestamp threshold,
announced before it can be crossed.\\
STATUS & A future live/off update for the named regions. It changes
selection eligibility, not existing voting weights.\\
MSIG FORCE & A target branch epoch, activation time and any authorized
status changes to apply with the Force. No retained height is supplied.\\
CLEAR & An explicit future removal of the named automatic cooldown;
ordinary ON does not perform this operation.\\\bottomrule
\end{tabularx}
\end{table}

\paragraph{Authority records.}
For each $(r,g)$, L1 supplies at most one immutable authority.
$\tagg{AUTH}_{r,0}$ fixes $\mathit{cfg}_r$, including the members,
weights, $F_i$, $\Delta_i$ and $K$ of each region, and fixes $\TB_r$.
It also fixes positive window lengths $W_{\rm start},W_{\rm run}$,
positive count targets $m_{\rm start},q_{\min}$, and cooldown
$C_{\rm auto}>0$. Counts are capped by the remaining DATA quota.
For $g>0$, only the leader schedule changes.
Nodes wait for required authority evidence. Later controls and slashing
never reweight an old certificate or reset a session quota.

Authority is indexed by $r$, not by a particular speculative branch.
A new epoch that reuses $r$ uses the same authority at that index.
Its session ID identifies a different engine. Rebinding an already
assigned authority is not permitted.

\paragraph{Control ingestion and activation.}
The complete L1 control prefix has a sequence number and a binding
history commitment. A verifier retains pending records as well as
already applied records. Ingestion does not imply activation.
Activation times need not follow sequence order: records becoming due
together are ordered by $(\theta_a,c)$. Malformed, unauthorized or
duplicate records do not grant a second action.

STATUS and CLEAR are independent controls. Configuration attached to
a Force applies only if that Force's target epoch is still current.
A stale Force is consumed without changing the new epoch. An independent
STATUS beside it still takes effect on its own terms.
Records are immutable after finalization. A revision is a new,
prospectively scheduled record, not an edit to past history.

\paragraph{Espresso positions and effective status.}
Let $J_q$ be Espresso block $q$, with authenticated timestamp
$\theta_q$. Positions also include namespace-transaction and record
indices. Heights and those indices give a total order even when
timestamps are equal. For a threshold $d$, define
\[
 q(d)=\min\{q:\theta_q\ge d\}.
\]
Its event is applied at the opening of $J_{q(d)}$, before that block's
EXPORTs. The prelude applies due independent status changes, then a
matching multisignature Force, then any automatic Force.

Write $\chi_E(q)$ for the eligibility snapshot after this prelude and
before the EXPORTs. It combines the L1 live/off map with automatic
cooldowns through that position. A region is eligible only if its
governance flag is on, its cooldown has expired, and its designated
authority has positive total weight. The historical snapshot is
authenticated by replay, not by an endpoint's statement.

OFF changes subsequent selections made from the affected snapshots.
It does not revoke an existing historical REGULAR witness. ON does not
clear cooldown or reopen an old epoch. A committee that needs to stop
the current branch uses a Force. No control changes an accepted DATA
prefix or retroactively changes a regional voting snapshot.

\paragraph{Genesis.}
The initial checkpoint fixes the chain identifier, cyclic order
$\mathcal L_R$, genesis DATA anchor, first replay position, Espresso
quorum trust anchor, control prefix and initial epoch $\beta_0$.
A synthetic boundary selects the first eligible region in
$\mathcal L_R$ at the first eligible authenticated position. It waits
if none exists. The initial authority and state are trusted or verified
from their specified sources; Espresso does not authenticate earlier
application history merely by starting this interpreter.
The genesis anchor includes synthetic commitment evidence and fixed,
authenticated L1 and Espresso references. Its pending boundary has
context $(\bot,-1,\bot)$ and opens numeric session 0. At its selection
height $q_s$, set $q^{\rm open}_0=q_s$ and
$\lambda^{\rm open}_0=\max\{\lambda(\text{genesis}),
\lambda(\tagg{AUTH}_{0,0})\}$, where the second term is the authority's
L1 publication height. A supplied later checkpoint must instead include
its current boundary, reference bounds and complete interpreter state.

\subsection{Stake and faults}
\label{sec:stake}
Member $v$ of region $i$ has voting weight $\omega_{i,v}>0$. Let
$W_i=\sum_{v\in V_i}\omega_{i,v}$. The stake configuration
fixes a bound $F_i<W_i/2$, and Byzantine members of region $i$ hold
total weight at most $F_i$. Every regional certificate needs
matching signatures from distinct members with total weight greater than
$F_i$. So every certificate has a correct signer, and the correct members
alone can form one. This replaces the $f+1$ of $n=2f+1$ in Part~I.

Membership and weights stay fixed for a whole session, including its
recoveries. Slashing and status updates do not change them. A later
session checks old certificates against the old snapshot. The bound
$F_i$ counts every key of the snapshot that is compromised while its
certificates are still accepted. Compromise of old keys after that is
outside this model.

Correct replicas keep their protocol state and committed blocks.
Signatures cannot be forged, and hashes are binding.



\subsection{Assumptions}
\label{sec:assumptions}
\begin{assumption}[Finalized L1 records]
\label{ass:authority}
At every finalized L1 height, the authority records and ordered control
history through that height are common and authenticated. Finalized
records never change. A verifier obtains a complete control prefix,
not an arbitrary subset of valid events.
\end{assumption}

\begin{assumption}[Authenticated synchronous communication and global clock]
\label{ass:network}
\leavevmode\par
\begin{itemize}
\item Messages use authenticated communication channels.
\item Within region $i$, the actual delay of a correct-to-correct
message is $\delta\le\Delta_i$; $\delta$ may differ across messages.
\item All nodes share a global clock.
\end{itemize}
\end{assumption}
The delay bound includes delivery and processing of complete validation
dependencies. Each replica records its own entry and voting events on
the global clock.

\begin{assumption}[Cross-region backbone]
\label{ass:backbone}
A backbone eventually delivers each session's committed blocks, with
their regional receipts, to all other nodes. Its delay has no fixed
bound.
\end{assumption}



\begin{assumption}[Honest multisignature committee]
\label{ass:committee}
The committee is one honest, coordinated control service. It submits
only authorized, immutable, future-dated STATUS, CLEAR and MSIG FORCE
records. It chooses activation times satisfying the advance-announcement
condition below. It keeps no private protected prefix and has no power
to select an older retained anchor.
\end{assumption}

\begin{assumption}[Bounded L1 finalization and visibility]
\label{ass:monitor}
A valid control operation actually submitted at real time $s$ is
finalized, available and verifiable by every correct online participant
by $s+T_{\rm L1}$. Participants include regional replicas, external
observers and certifiers. The bound includes inclusion, finalization,
delivery, retrieval and validation.
\end{assumption}

\begin{assumption}[Espresso history and time]
\label{ass:espresso}
Correct verifiers obtain compatible finalized Espresso prefixes and
authenticate every required namespace payload or its genuine absence.
Finalized history is retained and retrievable. Timestamps are
non-decreasing. If block $J_q$ first becomes final at real time
$t_E(q)$, then $\theta_q\le t_E(q)+\epsilon_E$, for the configured
maximum clock lead $\epsilon_E\ge0$.
\end{assumption}

\paragraph{Advance announcement.}
For a control submitted at $s$, the committee chooses
\[
 \theta_a>s+T_{\rm L1}+\epsilon_E+\eta,
 \qquad \eta>0.
\]
Thus the control is available to correct participants before its
first crossing can become final. This is a model condition, not a
conclusion from an average L1 or Espresso block time. The L1 service
bound already includes dissemination; no second announcement delay is
added. Offline or unsynchronized participants catch up before confirming
new output. Missing control history pauses verification.

\paragraph{Verification interfaces.}
Regional receipt validity and the application-validity checks below
are assumed sound under the stated regional model. A deployed
certifier must authenticate both source histories and bind the control
version used. A signature or TEE alone does not establish that an old
control snapshot is complete. The publication interface additionally
requires sound execution and DA correspondence checks.

\subsection{Conditions for progress}
\label{sec:progress}
The following conditions are separate from prefix safety.
\begin{itemize}
\item Correct leaders receive repeated turns that are actually entered,
including recovery initializers and DATA-producing turns across
sessions. Exposed equivocators are eventually excluded permanently
from later leader schedules.
\item Espresso continues to finalize retrievable history with growing
timestamps and includes this chain's valid exports.
\item Complete regional receipts and payloads reach at least one
working, funded publisher. Publisher retries and data transfer are
timely enough for some service interval to meet the configured windows.
\item Required authorities and recoverable native state become
available. A waiting boundary eventually has another eligible region.
The committee leaves working regions enough uninterrupted service time.
\item Correct clients' admissible requests reach the regions that
continue serving. Requests fit in DATA batches and have only finitely
many older requests ahead of them in each local queue.
\end{itemize}
An eventual backbone alone does not give a finite delivery deadline.
Window success therefore needs the additional timely-service condition;
otherwise repeated automatic Forces can prevent progress without
revoking any global DATA. Espresso failure also stops external time and
global progress. Switching the global finality service is outside this
protocol.

\section{Regional engine: definitions}
\label{sec:regional}
\Cref{tab:notation} lists the principal objects. Regional anchors and
view times retain their full session indices.
\begin{table}[htbp]\centering\small
\caption{Notation.}\label{tab:notation}
\begin{tabularx}{\linewidth}{@{}p{.32\linewidth}Y@{}}\toprule
Symbol & Meaning\\\midrule
$(r,g,t)$; $j$ & Session, recovery generation, leader term, and fresh position.\\
$\beta$; $\sigma_r$ & Branch epoch; full session identity.\\
$i_r$; $V_i$ & Session region; its regional members.\\
$\omega_{i,v},W_i,F_i$ & Member weight, total weight, and Byzantine bound $F_i<W_i/2$.\\
$\Delta_i$; $K$ & Regional synchronous delay bound; maximum DATA per term.\\
$\tagg{AUTH}_{r,g}$; $\mathit{cfg}_r$ & Immutable authority; fixed session snapshot.\\
$A_{r,g,t}$ & Anchor before the first fresh position of that view.\\
$h(B)$; $k_r(B)$ & DATA height; session physical position, including DRAIN.\\
$\TB_r,H_r,T_r$ & Session DATA quota, endpoint height, and terminal DATA.\\
$m_{r,g,t}$; $L(r,g,t)$ & Fresh term positions; designated leader.\\
$\QC_{r,g,t}(B)$; ANCHOR QC & Proposal certificate; opening re-certification.\\
$\FC_r(B)$ & Regional FINAL receipt, or its stated same-session coverage.\\
$Q_r^0,Q_v,Q^*$ & Inherited commitment evidence; entry lock; selected justification.\\
$e_v(r,g,t),\tau_v(P),a_v$ & Real-time entry, validated vote and quit events.\\
$\rho,\lambda(B)$ & Regional finalized L1 reference and its L1 height.\\
$\xi(B),\vartheta(B)$ & Signed Espresso selection reference and its timestamp.\\
$J_q,\theta_q$ & Espresso block and authenticated timestamp.\\
$\chi_E(q),\mathcal L_R$ & Historical pre-payload eligibility snapshot; region cycle.\\
$q_s,q_r^{\rm open}$ & Selection Espresso height; session reference lower bound.\\
$\lambda_r^{\rm open}$ & Separate L1 reference lower bound for a new session.\\
$G_u,R_u,C_{\rm pub}$ & Global DATA prefix; provisional branch; accepted publication prefix.\\
$T_{\rm L1},\epsilon_E$ & Full L1 announcement-service bound; maximum Espresso clock lead.\\
$W_{\rm start},m_{\rm start}$ & Startup-window length and DATA target.\\
$W_{\rm run},q_{\min},C_{\rm auto}$ & Running-window length, target and cooldown.\\
$E$; Ready & Regional entry proof; full boundary/state/authority validation.\\
\bottomrule\end{tabularx}
\end{table}

\subsection{Blocks, positions, and L1 references}
\label{sec:blocks}
The \defterm{anchor} $A_{r,g,t}$ is the block immediately before the first fresh
position of view $(r,g,t)$. We always write all three indices.
\begin{itemize}
\item $A_{r,0,0}$ is the DATA block inherited from the preceding
session; $A_{0,0,0}$ is genesis. It keeps the regional commitment evidence from the session that
produced it: a direct or covering receipt as defined in
\cref{sec:coverage}. Genesis has synthetic evidence.
\item For $g>0$, NEWVIEW chooses the \defterm{recovery anchor} $A_{r,g,0}$
(\cref{alg:recovery}). It is certified but may be uncommitted.
\item For $t>0$, $A_{r,g,t}$ is the terminal DATA block of term
$(r,g,t-1)$, certified by that term's terminal QC.
\end{itemize}
The anchor of a view never changes. The generation starts at
$A_{r,g,0}$; later terms use their own anchors $A_{r,g,t}$. ANCHOR voting
re-certifies only $A_{r,g,0}$, once per generation.

Each block has two positions.
\begin{itemize}
\item The \defterm{DATA height} $h(B)$ counts the DATA blocks on the chain
from genesis to $B$.
\item The \defterm{physical position} $k_r(B)$ counts DATA and DRAIN blocks
within session $r$, with $k_r(A_{r,0,0})=0$. This local count does not change
the signed fields of $A_{r,0,0}$.
\end{itemize}
A new block of session $r$ has the form
\[
 B=(r,k,h,\mathit{type},x,\Hash(B_{\rm parent}),\rho,\xi),
 \qquad \mathit{type}\in\{\tagg{DATA},\tagg{DRAIN}\},
\]
with $k=k_r(B_{\rm parent})+1$. A \defterm{DATA block} carries a batch $x$ of
requests and has $h=h(B_{\rm parent})+1$. A \defterm{DRAIN block} carries no batch
and has $h=h(B_{\rm parent})$.

\paragraph{L1 references.}
The \defterm{reference} $\rho$ names a finalized L1 block. It holds that L1
block's height $\lambda(B)$ and hash, and a registry proof for the L1
control records up to that height. These rules hold:
\begin{itemize}
\item along DATA ancestry, $\lambda$ never decreases;
\item the first DATA block of session $r$ has $\lambda\ge
\lambda^{\rm open}_r$, the L1 lower bound fixed by its boundary
(\cref{sec:witness}; genesis uses \cref{sec:l1});
\item a DRAIN block copies the reference of $T_r$.
\end{itemize}
The authenticated timestamp of the L1 block named by $\rho$ can be
read from its header. It is not a commit timestamp for $B$. Recovery
may commit $B$ much later (\cref{sec:late-receipt}). The
regional Vote step applies no receive-time freshness filter to $\rho$.

The L1 reference retains its regional validation role. It does not
choose the external Force time or the new Espresso selection snapshot.
Regional validity uses the fixed session authority. A later status update
does not change an old receipt's voting weights.

\paragraph{Espresso selection references.}
The additional outer field $\xi(B)$ names a finalized Espresso header
$J_q$ by height and hash, with its authenticated pre-payload eligibility
snapshot. Write $\vartheta(B)=\theta_q$. Its height does not decrease
along DATA ancestry, and the first DATA of session $r$ has
$q\ge q_r^{\rm open}$. DRAIN copies the selection reference of $T_r$,
just as it copies its L1 reference. These references and the session
context are bound by the proposal hash and signatures. They do not
change the native mini-block encoding.

A verifier authenticates the referenced history, directly or from its
validated cache. Global acceptance additionally requires the reference
to precede the containing EXPORT's Espresso block. No reference can
justify itself through that same EXPORT. Historical references carry no
receive-time freshness claim. The reference of $T_r$ starts the normal
successor search in \cref{sec:witness}.

\subsection{Session quota, term positions, and DRAIN}
\label{sec:quota}
\paragraph{1. Fix the session endpoint.}
Session $r$ may add $\TB_r\ge1$ DATA blocks after $A_{r,0,0}$. Its \defterm{last
DATA height} is
\[
 H_r=h(A_{r,0,0})+\TB_r.
\]
The DATA block at this height is $T_r$. Recovery never changes $H_r$.
The quota counts the DATA blocks after $A_{r,0,0}$ in a proposal's
ancestry, whether committed or not. DRAIN, resends, and ANCHOR votes
use no quota.

\paragraph{2. Produce DATA while space remains.}
If $h(A_{r,g,0})<H_r$, generation $(r,g)$ is a \defterm{DATA generation}.
Each term begins at its anchor $A_{r,g,t}$ and has
\[
 m_{r,g,t}=\min\{K,\;H_r-h(A_{r,g,t})\}
\]
new positions, numbered $j=0,\ldots,m_{r,g,t}-1$. At position $j$, the
DATA block has
\[
 h=h(A_{r,g,t})+j+1,\qquad
 k=k_r(A_{r,g,t})+j+1.
\]
A new position is called a \defterm{fresh position}; its block is a
\defterm{fresh block}. The anchor itself is not fresh. The last position
$j=m_{r,g,t}-1$ is the \defterm{terminal position}, and its QC is the
\defterm{terminal QC}.

If the terminal DATA height is below $H_r$, its QC admits term
$(r,g,t+1)$, with that terminal block as $A_{r,g,t+1}$. If it is $H_r$,
there is no next term. Only the last DATA term can have fewer than $K$
positions. The initial generation is always a DATA generation, since
$\TB_r\ge1$.

\paragraph{3. Use DRAIN when recovery starts at the endpoint.}
If $g>0$ and $h(A_{r,g,0})=H_r$, generation $(r,g)$ is a
\defterm{drain generation}. It has only term $t=0$ and position $j=0$,
with $m_{r,g,0}=1$. Its one fresh block is DRAIN, whose parent is
$A_{r,g,0}$ and whose justification is the generation's ANCHOR QC.
It carries no requests, keeps DATA height $H_r$, and has no next term.

DRAIN supplies a commit timer. ANCHOR voting starts no timer, so
re-certifying $A_{r,g,0}$ alone does not commit it. At height $H_r$, a
DATA child would exceed the quota. The DRAIN child's timer instead
commits its ancestors, including $T_r$.
\begin{itemize}
\item A DATA generation that reaches $H_r$ does not append DRAIN.
It waits for $T_r$'s timer or recovers after a quit.
\item A later drain generation may start from an earlier DRAIN.
Each interrupted recovery may add one DRAIN, so their total number
is not bounded.
\item The choice depends on the recovery anchor's DATA height, not on
whether a replica has already committed $T_r$.
\end{itemize}
DRAIN stays in the regional log. The external chain contains only DATA,
and the next session inherits $T_r$.

\paragraph{Example.}
Let $h(A_{r,0,0})=100$, $\TB_r=10$, and $K=4$, so $H_r=110$.
\begin{center}\small
\begin{tabular}{@{}lll@{}}\toprule
View & Anchor height & Fresh DATA heights\\\midrule
$(r,0,0)$ & 100 & 101--104\\
$(r,0,1)$ & 104 & 105--108\\
$(r,0,2)$ & 108 & 109--110\\\bottomrule
\end{tabular}
\end{center}
Recovery from DATA 107 starts a new generation whose first term produces
108--110. Recovery from certified DATA 110 starts a drain generation.
In both cases, the endpoint remains 110.

\subsection{Certificates and signed messages}
\label{sec:certificates}
\Cref{tab:certificates} lists the certificates of a session. Each needs
matching signatures from distinct members of $V_{i_r}$ with total weight
$>F_{i_r}$.

\begin{table}[htbp]\centering\small
\caption{Certificates of a session.}
\label{tab:certificates}
\begin{tabularx}{\linewidth}{@{}>{\raggedright\arraybackslash}p{.25\linewidth}>{\raggedright\arraybackslash}p{.24\linewidth}Y@{}}
\toprule
Certificate & Signed message & Purpose\\
\midrule
\defterm{Proposal QC} $\QC_{r,g,t}(B)$ & DATA-VOTE or DRAIN-VOTE & Justifies a
child. A terminal QC below $H_r$ also admits the next term.\\
\defterm{ANCHOR QC} & ANCHOR-VOTE & Re-certifies the opening $A_{r,g,0}$ in
generation $g$, before any fresh proposal.\\
\defterm{BLAME proof} & BLAME & Freezes a generation (\cref{sec:quit}).\\
\defterm{Regional receipt} $\FC_r(B)$ & FINAL & Proves that a correct replica
committed DATA block $B$.\\
\bottomrule
\end{tabularx}
\end{table}

A correct replica casts at most one fresh vote per $(r,g,t,j)$ and one
ANCHOR vote per $(r,g)$. DATA and DRAIN votes follow the same rule but
are different signed types. No QC and no ANCHOR vote starts a commit
timer.

\paragraph{Signed context.}
Every signature binds the session ID $\sigma_r$ and its branch epoch
$\beta$, the
voting configuration, session $r$, region $i_r$, and the message type.
This separates ordinary continuation from a later force fork
(\cref{sec:witness}). Messages of a view also bind the L1
authority and $(r,g,t)$. The other signed fields are:
\begin{itemize}
\item \textbf{PROPOSE and fresh votes:} the opening identity
(below), position $j$, block
type, physical position, DATA height, and block hash;
\item \textbf{NEWVIEW and ANCHOR-VOTE:} the identity of the chosen
opening. NEWVIEW also binds its justification and entry proof;
\item \textbf{STATUS:} the sender's entry lock and entry proof. No
opening is needed;
\item \textbf{BLAME:} the view only. It binds no opening, block hash, or
progress count. All complaints for a view thus sign the same statement,
even before NEWVIEW;
\item \textbf{FINAL:} session-stable fields only (below).
\end{itemize}
Message displays omit these context fields.

\paragraph{Rank and the base justification.}
Within session $r$, a QC formed in term $(r,g,t)$ for block $B$ has \defterm{rank}
$(g,t,k_r(B))$, compared lexicographically. The ANCHOR QC of generation
$g$ has rank $(g,0,k_r(A_{r,g,0}))$, just below the generation's first
fresh QC. A fixed total order on certified content breaks ties. Signer
sets and encodings do not matter. The \defterm{opening
identity} of generation $g$ is
$(g,k_r(A_{r,g,0}),h(A_{r,g,0}),\Hash(A_{r,g,0}))$. In this document,
``opening'' always means the generation's opening block $A_{r,g,0}$, as
in Part~I; a session is named by its session ID.

The anchor $A_{r,0,0}$ also has a \defterm{base justification} $Q_r^0$: the
anchor together with its direct or covering commitment evidence in the
boundary witness (\cref{sec:coverage,sec:witness}). Its justified block
is exactly $A_{r,0,0}$, even if the evidence uses a receipt for a higher
old-session block. That higher block is evidence only. $Q_r^0$ ranks
below every QC of session $r$.
Recovery uses it when no QC of the session is known. A QC of another
session is never a justification in session $r$, so it cannot replace
$A_{r,0,0}$.

\paragraph{Regional receipts.}
After a replica commits a DATA block $B$ of session $r$, by $B$'s own
timer or as an ancestor, it sends
\[
 \langle\tagg{FINAL},\mathit{cfg}_r,r,h(B),\Hash(B)\rangle_v
\]
to its region and to the backbone. Matching FINAL messages of weight
$>F_{i_r}$ form the \defterm{regional receipt} $\FC_r(B)$. The receipt carries $B$ and the ancestry needed to validate it, or
refers to a validated local cache.

\begin{definition}[\defterm{Valid regional receipt}]
\label{def:receipt}
A regional receipt $\FC_r(B)$ is valid when all of the following hold:
\begin{enumerate}[label=(\roman*)]
\item $B$ is a DATA block of session $r$, extends $A_{r,0,0}$, and has
$h(A_{r,0,0})<h(B)\le H_r$;
\item distinct members in session $r$'s fixed voting snapshot have
valid signatures on the same
$\langle\tagg{FINAL},\mathit{cfg}_r,r,h(B),\Hash(B)\rangle$,
with total weight $>F_{i_r}$ and the signed context specified above; and
\item $B$ and its required ancestry pass the block-validation rules of
\cref{sec:admission}.
\end{enumerate}
The verifier checks this evidence; it does not replay the signers'
timers. The weight threshold guarantees at least one correct signer,
who sends FINAL only after committing $B$.
\end{definition}
Here $\mathit{cfg}_r$ identifies session $r$'s \defterm{fixed voting snapshot}.
When we write $\FC(B)$ without a session subscript, we mean the receipt
from the session that produced $B$. The inherited anchor keeps that
receipt; it does not acquire a new one. A verifier outside
the region checks condition~(iii) from evidence, as
\cref{sec:admission} explains.
\begin{itemize}
\item FINAL binds no generation or term. So FINAL messages for one block
combine across recoveries. The block hash keeps the block's signed
context.
\item If a DRAIN commits $T_r$, replicas send FINAL for $T_r$, never for
the DRAIN.
\item Replicas of session $r$ send no FINAL for $A_{r,0,0}$. It keeps its original commitment evidence, which may come from an
earlier session if intermediate sessions added no DATA.
\item A replica may also commit $B$ and its uncommitted
ancestors when it obtains a valid $\FC_r(B)$ of its own session.
Completion does this for $T_r$ (\cref{alg:steady}). Such a commit
sends no FINAL, because the receipt already exists. Only Completion
also cancels timers. Receipts of other sessions serve only as history or anchor
evidence. They never change this session's locks or log.
\end{itemize}

\paragraph{Receipt coverage.}
\label{sec:coverage}
A valid $\FC_r(D)$ also proves commitment of every DATA ancestor $B$
of $D$ in the same session. At least one correct signer committed
$D$, which commits those ancestors. The \defterm{covering evidence} includes the
valid receipt for $D$ and the fully verified ancestry linking $D$ to $B$.
It does not claim that the signatures name $B$ directly.

This coverage never substitutes for a receipt from a different session.
It also does not replace the check that $B$ belongs to the authorized
external branch. A later Force may discard $D$ while keeping $B$; the
regional commitment evidence for $B$ still holds. External observers,
retained-anchor witnesses, and base justifications may use direct or
covering evidence. No extra FINAL messages or certificate round are
needed. We use ``\defterm{commitment evidence} for $B$'' for either form.

\subsection{Admission, validation, and engine states}
\label{sec:admission}
\paragraph{\defterm{Admission.}}
Replica $v$ admits views as follows.
\begin{itemize}
\item $(r,0,0)$ is admitted at $v$'s Ready event for the boundary that
opens session $r$ (\cref{def:ready}). The witness is the \defterm{entry proof}.
The NEWVIEW must choose exactly $A_{r,0,0}$.
\item $(r,g,0)$ with $g>0$ is admitted once $v$ holds
$\tagg{AUTH}_{r,g}$ and a quit proof for generation $g-1$. That proof is
the \defterm{entry proof}. The NEWVIEW is justified by a QC of an earlier
generation of session $r$, or by $Q_r^0$.
\item $(r,g,t)$ with $t>0$ is admitted once $v$ also holds the terminal
QCs of terms $0,\ldots,t-1$, each for DATA below $H_r$.
\end{itemize}
Term zero is admitted before its NEWVIEW arrives. Its progress checks
therefore run even if the initial leader is silent. Fresh proposals need
the generation's ANCHOR QC. This holds in every generation, including
generation $0$ of session $0$; unlike Part~I, genesis has no shortcut. A
message is \defterm{admissible} if its view is admitted and its signatures
are valid.

A verifier that is not a replica of the session has no
admission events. When it validates history (\cref{def:receipt},
\cref{sec:external}), a view counts as admitted if the evidence contains
its \defterm{admission witness}: the boundary witness for $(r,0,0)$, the chain of
entry proofs for $(r,g,0)$, and the earlier terminal QCs for $t>0$.

\paragraph{\defterm{Complete validation.}}
Validation checks signatures, authority, admission, block type, quota,
both positions, the L1 and Espresso reference rules, the parent rule
(\cref{tab:parent}), parent QCs, hash links, and application validity
throughout the required ancestry. Every relevant ancestor includes its
signed proposal and block data; a chain of certificate hashes alone is
not enough. Dependencies come from the message or a validated local
cache. Correct senders send \defterm{complete validation bundles}. A reference is
allowed only if the receiver can resolve it within the delay bound of
\cref{ass:network}. Bundles list each dependency once.

\paragraph{Evidence takes priority.}
Before validating a message or discarding a duplicate, inspect all signed
material in it for a quit proof (\cref{sec:quit}). A proof needs its own
valid signatures and admission witnesses. A valid proof counts even if
the message that carries it is invalid or incomplete. Compare incoming
material with all retained material, not only with objects inside the
message.

\paragraph{One proposal-processing rule.}
Every admitted fresh proposal of the current generation that $v$ has not
voted for invokes the Vote step. This includes proposals in attached
histories, proposals of earlier terms, and proposals whose blocks are
already committed as ancestors. A proposal's identity is
$(r,g,t,j,\Hash(B))$. The signer set or encoding of its parent QC does
not change it. Deduplicate by this identity after evidence inspection.
Old-generation proposals are retained without new votes or timers.
Future-generation messages wait for entry.

\paragraph{Engine states.}
A replica runs one engine for each session ID of its region.
Engines of different sessions remain separate, even if they use the
same session index. The engine of
session $r$ is in one of three states.
\begin{itemize}
\item \textbf{\defterm{Active}:} the current generation is admitted and unfrozen.
The engine proposes, votes, and changes terms.
\item \textbf{\defterm{Frozen}:} the current generation has a quit proof. The
engine stops proposals, votes, and term changes. It still processes
evidence, old votes, QCs, receipts, and recovery.
\item \textbf{\defterm{Complete}:} the replica holds $\FC_r(T_r)$ satisfying
\cref{def:receipt}.
Production in session $r$ is over, even if the external handoff is still
pending. The engine only serves data.
\end{itemize}
Committing $T_r$ locally does not complete the engine. Other members may
still need this replica's recovery votes before a receipt forms. Each
message goes to the engine identified by its session ID, so
evidence from one session or branch never freezes another. These engine
states do not describe external finality. The external chain may move
to another branch while an old engine still runs. L1 status
updates and Forces never change an engine's state, timers, or log; they
act only on the external branch (\cref{alg:rotation}).

\begin{table}[htbp]\centering\small
\caption{Parent selection for the next fresh position $j$ of term $(r,g,t)$.}
\label{tab:parent}
\begin{tabularx}{\linewidth}{@{}l Y Y@{}}
\toprule
Condition & $B_{\mathrm{parent}}$ & $Q_{\mathrm{parent}}$\\
\midrule
$j=0,t=0$ (DATA or DRAIN) & Opening $A_{r,g,0}$ & Its ANCHOR QC in $(r,g,0)$\\
$j=0,t>0$ & Anchor $A_{r,g,t}$, the terminal block of $(r,g,t-1)$ & Its proposal QC in $(r,g,t-1)$\\
$1\le j<m_{r,g,t}$ & Position $j-1$ of $(r,g,t)$ & Its proposal QC in $(r,g,t)$\\
\bottomrule
\end{tabularx}
\end{table}

\section{Regional protocols}
\label{sec:protocols}
Protocols 1 and 2 run inside the engine of one session $r$. They are
Protocols 1 and 2 of Part~I~\cite{part1} with the changes listed below.
\RR{Red text} marks the retained regional changes from Part~I.
The new outer selection reference is defined in \cref{sec:blocks}.
External replay and Forces do not change these regional timers.

\paragraph{Notation.}
All anchors and view times use their full session indices. The displayed
regional messages omit $r$, but their signatures bind the session ID
$\sigma_r$ and $r$. In these two protocols, $\Delta$ and $F$ mean
$\Delta_{i_r}$ and $F_{i_r}$.

\paragraph{Changes from Part~I.}
\begin{enumerate}[label=\arabic*.]
\item \emph{Stake.} Certificates need weight $>F$, not $f+1$ signatures.
\item \emph{Blocks.} A block also carries its session, physical
position, type, L1 reference $\rho$, and Espresso reference $\xi$
(\cref{sec:blocks}). Validity
checks all of them.
\item \emph{Quota.} A term has $m_{r,g,t}\le K$ positions. No term follows
$H_r$.
\item \emph{DRAIN.} A drain generation proposes one DRAIN. Its vote
counts as progress.
\item \emph{Receipts.} Replicas send FINAL after committing DATA. A
receipt for $T_r$ completes the engine (new step~5 of \cref{alg:steady}).
\item \emph{Timers.} Completion also cancels timers. A force never does.
\item \emph{Recovery entry.} After the $\Delta$ wait, a replica enters
when $\tagg{AUTH}_{r,g+1}$ is finalized. $Q_r^0$ can serve as the entry
lock and as the justification.
\item \emph{Genesis.} Session $0$ also certifies its anchor with an
ANCHOR QC before its first fresh proposal.
\item \emph{Sessions.} Each message goes to the engine of its session.
Session $0$ opens by bootstrap (\cref{alg:bootstrap}).
\end{enumerate}

\paragraph{Message-processing order.}\leavevmode
\RR{Replica $v$ first passes each message to the engine of the session
ID that the message binds. L1 updates go to \cref{alg:force}. Boundary witnesses and receipts of
other sessions go to \cref{alg:rotation,alg:external}.} The engine handles each received
message in this order:
\begin{enumerate}[label=(\arabic*)]
\item check all signed material for a quit proof;
\item check completeness and validity;
\item process attached DATA \RR{and DRAIN} proposals under Vote, parents
first;
\item process attached terminal QCs under step~4 of \cref{alg:steady},
in increasing term order;
\item process the message itself.
\end{enumerate}
When a message and a timer fall at the same time, $v$ handles the
message first.

\subsection{Steady state}
\label{sec:steady}
\Cref{alg:steady} gives the steady state of a session: proposals, votes,
timers, term rotation, and completion.

\begin{protocol}{alg:steady}{Propose, vote, commit, normal rotation\RR{, and completion}}
Replica $v$ runs these steps \RR{in the engine of session $r$}. Its
current term is $(r,g,t)$.
\begin{enumerate}[label=\textbf{\arabic*)},itemsep=3pt]
\item \textbf{Propose.} Run this step only if $v=L(r,g,t)$, $g$ is
unfrozen\RR{, and session $r$ is not complete}.
\RR{In a DATA generation,} for each position
$j=0,\ldots,\RR{m_{r,g,t}}-1$, in order:
 \begin{itemize}
 \item Wait until $v$ holds $B_{\mathrm{parent}}$ and
       $Q_{\mathrm{parent}}$ for $j$ (\cref{tab:parent}).
 \item Build \RR{$B=(r,k,h,\tagg{DATA},x,\Hash(B_{\mathrm{parent}}),\rho,\xi)$,
       with $k$ and $h$ from \cref{sec:quota} and the references
$\rho,\xi$ from \cref{sec:blocks}}.
       The batch $x$ follows \cref{sec:requests}.
 \item Broadcast
       $P=\langle\tagg{PROPOSE},g,t,j,B,Q_{\mathrm{parent}}\rangle_v$
       with its validation bundle \RR{to regional replicas and the backbone. }
 \item Process $P$ under Vote.
 \end{itemize}
 \RR{In a drain generation, propose once, at $(t,j)=(0,0)$. The block is
 a DRAIN with parent $A_{r,g,0}$, no batch, DATA height $H_r$, and the L1
 and Espresso references of $T_r$. Broadcast and process it as above.}
 Do not wait for the parent to commit. Never sign two different
 blocks for one position. A resend repeats the same $P$.

\item \textbf{Vote.} Run this step on every DATA \RR{or DRAIN} proposal
$P$ of term $(r,g_P,t_P)$ at position $j_P$. Here $t_P$ may be earlier
than $t$. The step also applies to proposals inside attached histories.
 \begin{enumerate}[label=(\alph*)]
 \item If $P$ or its attached material contains a quit proof for $g$,
       run Quit and stop.
 \item If $g_P<g$, retain $P$ and stop. If $g_P>g$, buffer $P$
       until $v$ enters generation $g_P$.
 \item If $g$ is frozen \RR{or session $r$ is complete}, retain $P$ and
       stop.
 \item If a dependency is missing, retain $P$, request the dependency,
       and stop. Rerun this step when it arrives.
 \item If $P$ is invalid, discard it and stop. \RR{Validity includes the
       block type, the quota, both positions, and the L1 and Espresso references.}
 \item If $v$ has already voted for $P$, stop.
       A duplicate never restarts a timer.
 \item Otherwise, let $\tau_v(P)$ be the current time. Forward $P$ and
       its bundle to all \RR{regional replicas and the backbone. }
       Broadcast
       $\langle\tagg{DATA\mbox{-}VOTE},g,t_P,j_P,\Hash(B)\rangle_v$
       \RR{(DRAIN-VOTE if $B$ is a DRAIN)}.
       Start the commit timer of $P$. It expires at $\tau_v(P)+2\Delta$.
 \end{enumerate}

\item \textbf{Commit.} When the timer of $P$ expires, commit $B$ and
its uncommitted ancestors. \RR{For each DATA block $D$ of session $r$
that this commits, send
$\langle\tagg{FINAL},\mathit{cfg}_r,r,h(D),\Hash(D)\rangle_v$ to the region
and the backbone. A committed DRAIN yields FINAL for $T_r$, never for
itself. Relay each receipt of session $r$ that $v$ obtains to the region
and the backbone.} Only Quit \RR{and Completion} cancel a timer. Entering a
new term does not cancel it. A QC never starts a timer. \RR{A force
instruction never cancels a timer.}

\item \textbf{Terminal QC and rotation.} Run this step the first time
$v$ obtains a valid terminal QC for a term $(r,g,t')$.
``Obtains'' means that $v$ assembles the QC from votes or finds it in
any message, including as another proposal's $Q_{\mathrm{parent}}$.
 \begin{enumerate}[label=(\alph*)]
 \item Process each unseen proposal in the QC's history under Vote,
       parents first.
 \item If $g$ is now frozen, stop. \RR{If the QC certifies DATA at
       $H_r$, or a DRAIN, stop: no next term exists.}
 \item Send the QC and its validation bundle to $L(r,g,t'+1)$.
 \item If $t\le t'$, enter $(r,g,t'+1)$ and set $e_v(r,g,t'+1)$ to the
       current time. Cancel the progress checks of all earlier
       terms.
 \end{enumerate}
 A QC never moves $v$ to an earlier term. It never changes a recorded
 entry time.

\item \RR{\textbf{Completion.} Run this step the first time $v$ obtains a
regional receipt $\FC_r(T_r)$ satisfying \cref{def:receipt}.}
 \begin{itemize}\color{rrRed}
 \item Commit $T_r$ and its uncommitted ancestors.
 \item Mark session $r$ complete. Stop its proposals, votes, progress
       checks, term changes, and recovery. Cancel its unexpired timers.
 \item Relay $\FC_r(T_r)$ to the backbone.
 \item Keep serving data. \Cref{alg:rotation} runs the handoff.
 \end{itemize}
 \RR{A force instruction never runs this step.}
\end{enumerate}
\end{protocol}

\Needspace{9\baselineskip}
The \defterm{acceptance time} $\tau_v(P)$ is when complete validation and the first
vote occur, not an earlier arrival of an incomplete packet. Each replica
records its own entry time for each term. In session $r$, $e_v(r,0,0)$ is
set at Ready, in step~1 of \cref{alg:bootstrap}. For $g>0$, $e_v(r,g,0)$ is
set in step~3 of \cref{alg:recovery}. For $t>0$, $e_v(r,g,t)$ is set in
step~4 of \cref{alg:steady}. A term's progress checks therefore start
when the term is admitted locally, not when its leader's first proposal
arrives. A silent new leader cannot delay them.

A proposal or QC at height $H_r$ does not end the session. Without a
quit, the timer of $T_r$ commits it; no later block is needed. Only a
receipt for $T_r$ completes the engine.

\subsection{Quit proofs and recovery}
\label{sec:quit}
A quit proof freezes a generation. \Cref{alg:recovery} then moves the
session to the next generation.

\begin{protocol}{alg:recovery}{Blame, quit, and recovery}
Replica $v$ runs these steps \RR{in the engine of session $r$}. Its
current term is $(r,g,t)$.
\begin{enumerate}[label=\textbf{\arabic*)},itemsep=3pt]
\item \textbf{Blame.} \RR{Let $q=m_{r,g,t}$ in a DATA generation and
$q=1$ in a drain generation.} For each $p=1,\ldots,\RR{q}$, run a
progress check at time
\[
 e_v(r,g,t)+(2p+4)\Delta.
\]
\RR{Schedule the check for $p=1$ at entry, even if the opening is not yet
known. Schedule the others once $q$ is known. A check whose time has
already passed runs at once.}
At that time, broadcast $\langle\tagg{BLAME},g,t\rangle_v$ if all of
these hold:
 \begin{itemize}
 \item $v$'s current term is still $(r,g,t)$, $g$ is unfrozen\RR{, and
       session $r$ is not complete};
 \item $v$ has cast DATA \RR{or DRAIN} votes for fewer than $p$ positions
       of term $(r,g,t)$ (ANCHOR votes do not count);
 \item $v$ has not yet sent BLAME for $(r,g,t)$.
 \end{itemize}

\item \textbf{Quit.} Run this step the first time $v$ holds a quit proof
for its current generation $g$.
 \begin{itemize}
 \item Freeze $g$. Stop all progress checks, proposals, votes, and
       term changes in $g$. Cancel every unexpired commit timer of $g$.
 \item Set $a_v$ to the current time. Broadcast the proof.
 \item Keep all commitments. Keep collecting votes and QCs of $g$.
       They can raise the entry lock chosen in step~3. \RR{Keep
       processing regional receipts (step~5 of \cref{alg:steady}).}
 \item If the proof is of type (b) or (c), submit it to L1. Do
       not wait for L1 to process it.
 \end{itemize}

\item \textbf{Status and entry.} \RR{Skip this step if session $r$ is
complete.} \RR{Wait until at least $a_v+\Delta$ and until
$\tagg{AUTH}_{r,g+1}$ is finalized. Set $e_v(r,g+1,0)$ to the time both
conditions hold, and read that authority from L1.} Then:
 \begin{itemize}
 \item Set the entry lock $Q_v$ to the highest-ranked QC
       (DATA\RR{, DRAIN,} or ANCHOR) that $v$ knows from generations
       $0,\ldots,g$\RR{, or to $Q_r^0$ if it knows none}.
       $Q_v$ never changes after this. QCs that arrive later do not
       update it.
 \item Enter $(r,g+1,0)$ as unfrozen and start its progress checks.
       \RR{Keep $A_{r,0,0}$, $H_r$, and the voting snapshot of session $r$.}
 \item Send $\langle\tagg{STATUS},g+1,0,Q_v,E\rangle_v$ and its
       validation bundle to $L(r,g+1,0)$. Here $E$ is the quit proof for
       $g$, which is also the entry proof of $g+1$.
 \item Process any buffered proofs for $g+1$.
 \end{itemize}

\item \textbf{New-view.} Run this step only at $L(r,g+1,0)$. Wait $2\Delta$
after its own entry. If $g+1$ is still current and unfrozen:
 \begin{itemize}
 \item Let $Q^*$ be the highest-ranked QC from generations $0,\ldots,g$
       that the leader knows, including the QCs in STATUS
       messages\RR{; if it knows none, let $Q^*=Q_r^0$}.
 \item Let $A_{r,g+1,0}$ be the block that $Q^*$ certifies. \RR{If
       $h(A_{r,g+1,0})<H_r$, generation $g+1$ is a DATA generation.
       Otherwise it is a drain generation.}
 \item Broadcast
       $\langle\tagg{NEWVIEW},g+1,0,A_{r,g+1,0},Q^*,E\rangle_{L(r,g+1,0)}$
       with its validation bundle. Choose only one opening block for
       $g+1$.
 \end{itemize}

\item \textbf{Anchor vote.} On receiving a NEWVIEW for $g+1$:
 \begin{enumerate}[label=(\alph*)]
 \item If it contains a quit proof for $g+1$, run Quit and stop.
 \item Validate its leader, entry proof $E$, opening block $A_{r,g+1,0}$,
       $Q^*$, and dependencies. If data are missing, request them and
       wait without voting. If the message is invalid, discard it.
 \item Vote only if all of these hold: $g+1$ is current and unfrozen; $v$
       has cast no ANCHOR vote in $g+1$; and $\rank(Q^*)\ge\rank(Q_v)$.
 \item To vote, forward the NEWVIEW to all replicas. Broadcast
       \[\langle\tagg{ANCHOR\mbox{-}VOTE},g+1,0,\Hash(A_{r,g+1,0})\rangle_v.\]
       Start no timer. Use no fresh position.
 \end{enumerate}
 Matching ANCHOR votes \RR{of total weight $>F$} form the ANCHOR QC.
 When $L(r,g+1,0)$ obtains it, the leader starts Propose in
 \cref{alg:steady} at $j=0$\RR{; in a drain generation, this is the
 DRAIN}. Every replica, including one that cast no ANCHOR vote,
 processes DATA \RR{and DRAIN} proposals backed by a valid ANCHOR QC
 under Vote.
\end{enumerate}
\end{protocol}

\paragraph{\defterm{Complete quit proofs.}}
For an admitted generation $g$ \RR{of session $r$}, a proof consists of
either:
\begin{enumerate}[label=(\alph*)]
\item distinct signatures \RR{of total weight $>F$} on the same
$\langle\tagg{BLAME},g,t\rangle$ for an admitted term, with that term's
admission witness;
\item two different admissible DATA \RR{or DRAIN} blocks signed by
$L(r,g,t)$ at the same $(r,g,t,j)$, with their admission \RR{and quota}
witnesses; or
\item two valid NEWVIEWs signed by $L(r,g,0)$ selecting different opening
blocks for generation $g$.
\end{enumerate}
\RR{A DATA block beyond the quota, or a DRAIN outside a drain generation,
is invalid. It is rejected and never counts for (b).}
Different openings count even when their blocks are comparable, because
they fix different starting heights $h(A_{r,g,0})$. The height formula then
maps the same position to different heights.

Different QCs for the same opening block are not equivocation. Evidence
from any admitted term can freeze its generation, including a term not
yet locally reached. Old-generation proofs are archival and may be
submitted to L1. \RR{Evidence from another session never freezes this
engine.} Future-generation proofs are buffered until entry. Neither
resets an unrelated recovery.

Publishers consume the same FINAL and receipt broadcasts; no publisher
acknowledgement is required for regional commitment or Completion.

The entry-lock test governs ANCHOR voting only. A replica that has
committed $T_r$ locally still takes part in recovery until it holds
$\FC_r(T_r)$, because other members may need its votes.

\section{Regional and global finality}
\label{sec:external}
Regional commitment and global acceptance are separate events.
A node follows a production branch $R_u$ and a continuous global DATA
prefix $G_u$. The branch may extend several normal sessions beyond
$G_u$. A Force closes that extension, never $G_u$ itself.

\subsection{Regional finality}
\label{sec:regional-finality}
\begin{definition}[\defterm{Regional finality}]
\label{def:regional-finality}
A DATA block $B$ has regional finality on a specified production branch
when its boundary history assigns it to its claimed region and session,
it has valid direct or same-session covering commitment evidence, and
the DATA path to it is valid and continuous. Each session on that path
supplies its own authority and commitment evidence.
\end{definition}
The region is derived from the boundary before its receipt is checked.
A valid receipt does not authorize another region's DATA or a different
session. Current live/off state does not replace historical authority.
When a Force closes an unconfirmed suffix, the receipts remain evidence
about the old regional engine. They no longer authorize that suffix to
extend the global prefix.

\paragraph{Remark (Commitment by one replica and by all replicas).}
Here ``all replicas'' means all correct voters of the session that
produced $B$. Consider that session's regional engine, with recovery
authorities set in advance.
\begin{enumerate}[label=(\arabic*),itemsep=4pt]
\item \textbf{What a receipt proves.} A valid direct or covering
receipt proves that at least one correct replica has committed $B$.
Other correct replicas may still have running timers, or may have
canceled their timers on Quit. Regional finality does not assert that
they have all committed already.

\item \textbf{When all replicas commit.} Regional safety preserves
$B$ through recovery: the selected anchor is $B$ or a descendant of
$B$. ANCHOR voting re-certifies that block but starts no commit timer.
Suppose a correct replica votes for and forwards a fresh DATA or DRAIN
proposal for a descendant $D$ at time $s$. If no correct replica freezes
that generation through $s+3\Delta_{i_r}$, all correct replicas commit
$B$ by that time. Replicas that have not completed the session vote by
$s+\Delta_{i_r}$; their timers commit $D$ and its ancestors at most
$2\Delta_{i_r}$ later. Completion may cancel a timer, but first commits
$T_r$ and its ancestors, including $B$. In particular, a DRAIN need not
itself commit at a replica that completes through the receipt for $T_r$.

\item \textbf{Why one honest leader is not enough.} An honest leader
may finish its term before its timers expire. A Byzantine successor
can then equivocate and freeze the generation. What is needed is the
uninterrupted commit window in (2). Under regional safety and the
progress conditions of \cref{sec:progress}, including permanent
equivocation exclusion and fair leader schedules, the engine eventually
obtains such a window while it continues to run, or completes through
the terminal receipt. Earlier recoveries do not revoke $B$. These
conditions do not give a fixed bound on when the window begins.

\item \textbf{When a replica can know.} A correct regional replica
that directly commits $D$ by its timer at time $c$ can conclude that
all correct regional replicas have committed $B$ if it remains Active
in the same generation through $c+2\Delta_{i_r}$, after processing
messages due at that time. Other replicas' timers are due by
$c+\Delta_{i_r}$. A Quit that canceled one of those timers would reach
this replica by $c+2\Delta_{i_r}$ and freeze its generation. A replica
whose timer was canceled by Completion has already committed $B$.
Normal term rotation does not affect this test.

\item \textbf{What an external observer can know.} A receipt alone
does not establish that uninterrupted window. The backbone has no
fixed delay bound, so receiving no quit proof is not evidence that no
replica quit. Thus receipt arrival starts no fixed wait that certifies
commitment by all regional replicas. Espresso global acceptance protects the external prefix against a Force.
It does not certify that every regional replica has completed its local
commitment.
\end{enumerate}



\subsection{Published records and the global interpreter}
\label{sec:exports}
A \defterm{publisher} carries regional results to Espresso. Any party
may perform this role. A publisher needs the full data and evidence;
it cannot replace the regional receipt with its own signature.
Replicas relay FINAL messages and assembled receipts over the backbone
as before. Publishers obtain those bundles, submit them, and retry
unaccepted ranges from the latest verified global base.

An \tagg{EXPORT} record binds the chain identifier, branch epoch,
expected global base, an ordered nonempty interval of DATA, and the
evidence needed to validate that interval. Evidence includes regional
receipts, ancestry, REGULAR boundaries, historical authorities and
native application context. A direct or covering receipt is supplied
for each session's DATA. A receipt from a later session does not cover
an earlier session. DRAIN may occur in supporting regional ancestry,
but is never exported as application DATA.

An EXPORT can contain the missing end of region A and the beginning
of region B. It therefore need not wait for a separate L1 publication
or for each ancestor to be submitted in its own Espresso transaction.
The entire record is checked on a temporary state. Either all of its
DATA, counts and normal boundary transitions are accepted, or none are.
Another receipt encoding or signer set does not create another DATA
position.

\paragraph{Acceptance.}
After the block's prelude in \cref{alg:force}, a record is accepted
only if all of these checks succeed:
\begin{enumerate}[label=(\roman*)]
\item the expected base equals the current $G$, including its prefix
commitment and session provenance;
\item the record's epoch is open, and every block is assigned to the
session reached by the validated normal boundary chain;
\item the DATA heights begin at $h(G)+1$, parents and hashes agree,
and no input or intermediate session is omitted or repeated;
\item regional commitment evidence, application validity and quota
checks succeed for every block;
\item each Espresso selection reference used by the record lies
strictly before its containing Espresso block, and its historical
eligibility snapshot is authenticated;
\item an epoch created by a Force in $J_q$ first exports in a block
strictly after $J_q$.
\end{enumerate}
The record is then committed atomically in namespace order.
Each new DATA increases the appropriate export-session count. On its
terminal DATA, that session's export window is cancelled. Its normal
successor becomes the next export session when the fixed selection
rule resolves. Its regional engine may have started much earlier.

A valid terminal DATA may enter $G$ even if no successor is yet
eligible. The interpreter keeps that terminal anchor and enters
\defterm{Waiting}. It runs no DATA window while no export duty is
assigned. A proven absence of a candidate is different from unavailable
evidence: missing history or authority evidence pauses the relevant
verification and is never treated as an empty candidate set.
When deferred evidence is obtained, replay resumes at the same canonical
position; local retrieval time never supplies a new public timestamp.

Malformed records and well-formed but invalid candidates have no
effect on $G$ or its counts. A future-height record skipped before its
ancestors are accepted must be submitted again. The interpreter never
goes back to make an earlier rejected record global. An incomplete
namespace download is not a malformed record: it prevents the block
from being processed at all.

\subsection{Global finality}
\label{sec:global-finality}
Let $S_u$ be observer $u$'s verified interpreter state. It contains $G_u$,
the current epoch and export duty, windows, eligibility state, pending
controls, and source-verification checkpoints.
Initially it is a trusted checkpoint or the state derived from genesis.

\begin{definition}[\defterm{Global finality at an observer}]
\label{def:global-finality}
Observer $u$ treats DATA block $B$ as globally final when it has
authenticated and replayed a finalized Espresso prefix containing an
accepted EXPORT that adds $B$ to $G_u$. Replay uses complete applicable
L1 control history and the event order of \cref{alg:force}. All DATA
predecessors are already accepted or are accepted in the same EXPORT.
\end{definition}

No extra waiting period follows this verification. A later Force
preserves the entire existing $G_u$, including its original block and
session identities. An old-epoch record after closure can still be
finalized as raw Espresso bytes, but it is not an accepted chain result.
Espresso inclusion alone, or a receipt alone, does not establish this
definition.

This is input-prefix finality. Espresso does not execute the native
transactions or prove their resulting state. The regional validity
contract and the execution-verification interface remain responsible
for that property. Finality of a state claim or withdrawal on L1 is a
separate stage. Publication lag never permits skipping globally final
input.

\subsection{References do not provide a commit clock}
\label{sec:late-receipt}
A block's L1 reference and Espresso selection reference are fixed when
it is proposed. A later recovery or DRAIN may commit it long afterwards.
Neither old reference starts an export window or establishes global
finality. An export window starts when replay assigns the next DATA
duty, and counts only newly accepted DATA from finalized Espresso
records. Regional timers continue to use the real-time clock of
\cref{ass:network}, not Espresso timestamps.

\section{Cross-region rotation}
\label{sec:xprotocols}
\subsection{Anchors, successor selection, and boundary identity}
\label{sec:witness}
Let $\tagg{NextLive}(i,M)$ be the first other eligible region after $i$
in the fixed cyclic list $\mathcal L_R$. Search at most one cycle.
If no such region exists, return $\bot$. Eligibility uses the historical
governance flag and cooldown in $M$, together with the designated
successor authority. It never changes an existing session's voters.
For the synthetic genesis source $\bot$, search the entire list from
its first region instead.

A pending boundary stores its source region $i_b$, predecessor index
$r_b$ and predecessor session ID $\sigma_b$. Its next index is $r_b+1$.
A regular terminal boundary stores the completed session's context.
A Force with an assigned export duty captures that duty's context.
A Force during Waiting reuses the stored context, replaces the pending
boundary and searches again from its triggering position. It does not
increment an unopened session. This rule also covers repeated Forces
before any successor opens.

A regular handoff inherits $A_{r+1,0,0}=T_r$.
A Force inherits exactly the current global anchor $G$ at the opening
of its triggering Espresso block. Its successor index is $r_c+1$,
where $r_c$ is the predecessor index in the captured context. With an
assigned duty, this is its export-session index; during Waiting, it is
the stored $r_b$. When the export session has accepted no DATA, $G$ may
belong to an earlier session. It keeps that session's original receipt and provenance.

\paragraph{Selection reference.}
For a regular boundary, let $q_0$ be the height in $\xi(T_r)$.
For a forced boundary, let $q_0$ be the triggering Espresso height.
The source region for the cyclic search is $i_r$ for a normal handoff,
and the source region in the captured context for a Force. It is not inferred from
the producer of $G$ or from a speculative successor.

Given authenticated Espresso history through $\ell$, define
\[
 \mathcal S(\ell)=\{q:q_0\le q\le\ell,\quad
       \tagg{NextLive}(i_{\rm source},\chi_E(q))\ne\bot\}.
\]
The candidate regions are those of the designated next authority.
If this set is empty, the boundary waits. Otherwise its selection
height and successor are
\[
 q_s=\min\mathcal S(\ell),\qquad
 i_{\rm next}=\tagg{NextLive}(i_{\rm source},\chi_E(q_s)).
\]
An authenticated prelude snapshot fixes each value of $\chi_E(q)$.
Extending the history can only add later candidate positions, so it
cannot change an already established minimum. Missing evidence is not
evidence that a candidate is absent. Nodes with insufficient history
wait for it.

For a normal successor, $q_s$ is available from existing Espresso
history; $T_r$ itself need not yet be published there. A later OFF does
not reselect this boundary. The reference may be historical: this
version does not claim a receive-time freshness guarantee. Immediate
closure of such an already authorized branch requires a Force.

\paragraph{Boundary witnesses.}
A witness contains the mode REGULAR, AUTO or MSIG, predecessor session
and epoch, source context, inherited anchor with original commitment evidence,
successor and authority, $q_s$, and source proofs. REGULAR also proves
the predecessor's exact quota and terminal receipt. AUTO proves the
failed replay window. MSIG proves the authorized finalized L1 record
and its first-crossing position. A forced witness authenticates the
complete replay prefix needed to identify $G$, not just one included
block.

Set $q^{\rm open}_{\rm next}=q_s$. The separate L1 lower bound
$\lambda^{\rm open}_{\rm next}$ is the maximum of the anchor's L1
reference height and the L1 height that published the next initial
authority. These are different indices: the first constrains Espresso
selection references; the second constrains regional L1 references.

\paragraph{Epochs and session IDs.}
A normal handoff retains $\beta$. A Force at $J_q$ closes the current
epoch and creates a new epoch from the chain identifier, old epoch,
canonical cause ID, triggering header hash, and retained anchor identity.
The cause ID binds the control sequence for MSIG, or the failed window
and export-session ID for AUTO. It does not depend on a signer set.
The new epoch can exist in Waiting before a successor is selected.
Its identity therefore does not depend on that successor.

The session opened by a boundary has identity
\[
 \sigma_{\rm next}=\Hash\bigl(\beta_{\rm next},\sigma_{\rm prev},
   \text{mode},\operatorname{id}(A_{\rm next,0,0}),
   i_{\rm next},q_s,\text{authority ID},\text{cause ID}\bigr).
\]
Here $\sigma_{\rm prev}$ is the predecessor ID stored in the boundary
context, including when a Force replaces a Waiting boundary.
REGULAR uses a fixed null cause ID. Its fields are known at handoff;
the identity does not depend on the future EXPORT position of $T_r$.
An external DATA identity includes its session ID and block hash.
Numeric session and DATA heights alone do not identify an old branch.

\subsection{Regular and forced handoff}
\begin{protocol}{alg:rotation}{Track production and rotate regions}
Each node follows the branch rooted in its verified global history.
\begin{enumerate}[label=\textbf{\arabic*)},itemsep=3pt]
\item \textbf{Regular completion.} On a valid terminal receipt
$\FC_r(T_r)$, verify the regional source, the exact DATA quota and the
signed selection reference. If the epoch remains open, select the
successor by \cref{sec:witness}, form the REGULAR witness, and send it
with full state and history to that successor. If no candidate is
available, retain the anchor and wait. Do not wait for global finality
or a new L1 publication. A future-dated Force does not block this path.

\item \textbf{Forced boundary.} On a verified AUTO or MSIG boundary
from \cref{alg:force}, preserve all of $G_u$ and close the old epoch
to further global output. Archive its unconfirmed production suffix,
including speculative normal successors. Keep their requests pending.
Fetch the retained anchor, its original evidence and execution state.
Select the successor or retain the boundary in Waiting. Send the full
forced witness to the selected members.

\item \textbf{Separate engines.} Route every message by its session
ID. New sessions use new locks, votes, entry times and protocol state.
Old engines keep their regional logs and timers. This external handler
does not run Quit or Completion for them. Closed-epoch outputs remain
archival even if their region later becomes eligible again.
\end{enumerate}
\end{protocol}

\subsection{Ready and the common bootstrap}
\label{sec:ready}
\begin{definition}[\defterm{Ready}]
\label{def:ready}
A selected member is Ready after completely validating its boundary,
inherited anchor, missing application state and matching finalized
authority. A REGULAR witness proves quota completion. A forced witness
also proves the finalized triggering Espresso block and the replay
state that closes the old epoch.
\end{definition}
A proposal QC, unfinalized Espresso header, pending L1 Force or live
status alone does not grant entry. Each member relays the complete
bundle at Ready. Duplicate witnesses do not reset entry state.

\begin{protocol}{alg:bootstrap}{Bootstrap the selected region}
Replica $v\in V_{i_{r+1}}$ runs these steps for one validated
session $\sigma_{r+1}$. Here $r$ is the predecessor index stored by
the boundary, namely $r_c$ for a forced boundary, and
$\Delta=\Delta_{i_{r+1}}$.
\begin{enumerate}[label=\textbf{\arabic*)},itemsep=3pt]
\item \textbf{Enter.} At Ready, install $A_{r+1,0,0}$,
$\tagg{AUTH}_{r+1,0}$, $\TB_{r+1}$, and the separate bounds $q^{\rm open}_{r+1}=q_s$ and
$\lambda^{\rm open}_{r+1}$ from the boundary. Set the entry lock
to $Q^0_{r+1}$. Enter $(r+1,0,0)$ unfrozen, record its entry time once,
and start progress checks. Send STATUS with this lock and boundary
witness to the initial leader. Process buffered proofs for this session.
\item \textbf{Choose the anchor.} The initial leader waits $2\Delta$
after its own entry. If the view remains current and unfrozen, it sends
NEWVIEW choosing exactly $A_{r+1,0,0}$, justified by $Q^0_{r+1}$ and
the boundary witness.
\item \textbf{Certify and produce.} Members apply the ANCHOR-vote step
of \cref{alg:recovery}. On the ANCHOR QC, the leader starts
\cref{alg:steady} at DATA height $h(A_{r+1,0,0})+1$ and physical
position 1. Later local failures use the same recovery protocol.
\end{enumerate}
\end{protocol}

Let $e=\min_v e_v(r+1,0,0)$ over correct members $v$, the
first correct Ready time. The complete-bundle relay makes
all correct members enter by $e+\Delta$. With a correct initializer and
no quit, NEWVIEW is sent by $e+3\Delta$, its ANCHOR QC reaches the
leader by $e+5\Delta$, and the first DATA proposal reaches all correct
members by $e+6\Delta$. This meets the first progress check. Transfer
before Ready remains subject to the backbone's eventual-delivery rule.



\paragraph{The public export clock.}
Ready starts regional timers at each replica, as above. It does not
start a public delivery window. Replay opens that window when it
assigns the next DATA duty: at the containing EXPORT's timestamp for
a normal successor already selectable, at the triggering block's
timestamp for an immediate forced successor, or at the first eligible
selection block after protocol Waiting. Local fetch time never supplies
the clock. Unavailable validation evidence pauses replay at its
canonical position rather than manufacturing a new opening time.

\section{The Force handler}
\label{sec:force}
Forces act on the external branch. They do not alter the regional
DATA/DRAIN engine. The interpreter derives an automatic Force itself
or applies a previously announced multisignature Force.
Both retain the same object: the entire global prefix immediately
before the triggering Espresso block's payload.

\subsection{Export duties and block-missing windows}
\label{sec:windows}
An export duty belongs to the session responsible for the next DATA
height after $G$. Its window parameters come from its fixed initial
authority, not from a later parameter update.
Let $R$ be its remaining DATA quota when a window opens. All targets
below are positive because completion removes the duty before $R=0$
can open another window.

At public opening timestamp $\theta_0$, create the startup window
\[
 [\theta_0,\theta_0+W_{\rm start}),\qquad
 m=\min\{m_{\rm start},R\}.
\]
After it succeeds, use consecutive fixed windows of length
$W_{\rm run}$, each with target $m=\min\{q_{\min},R\}$ measured from
that window's opening. The window stores its start, deadline, target,
accepted count and phase. A new block increments the count only after
its EXPORT is accepted into $G$. Each DATA is counted once in its own
session. Empty DATA counts under the same quota rule; DRAIN, heartbeats,
proposal QCs, duplicate records and future-session data behind a gap
do not count.

At the first Espresso block crossing the deadline, compare the count
with the target. Success opens the next fixed window and resets its
count. Extra output is not banked as unlimited credit for later windows.
If several deadlines have passed, settle the successful windows in
order and stop at the first failure. It produces one automatic Force.
A terminal DATA accepted before its deadline cancels that session's
remaining checks, even if the final window has not expired.

The intervals are half-open. DATA with a timestamp equal to the
deadline is too late for that window. The check runs before the first
block's payload at that timestamp, not after all blocks that may
eventually share the timestamp.

\paragraph{What a failure establishes.}
Failure means that the next required output was not delivered to the
global prefix at the specified rate. It does not identify a faulty
proposer, publisher, region or network. A backbone delay or Espresso
outage can cause the same observation. It is a recovery condition,
not automatically slashable equivocation or a proof of censorship.

\subsection{Scheduled control and two kinds of Force}
\label{sec:force-stages}
An MSIG Force binds the chain, target epoch $\beta$, activation
timestamp $\theta_a$, control sequence and authorized configuration
changes. Its identity is the hash of those canonical fields. It
contains no retained DATA height. The committee submits it under
\cref{ass:committee,ass:monitor} and serves its finalized evidence.
Signing alone does not start the submission bound.

L1 finalization makes the instruction known, but does not yet close
the branch. Its effect begins only at $J_{q(\theta_a)}$.
Until then, normal production, normal handoffs and valid exports can
continue. A node does not veto future global output merely because a
Force has been announced.

An automatic Force binds the chain, old epoch, export-session ID,
failed-window identity and triggering Espresso header. Its cause
deadline is the failed window's deadline; its realized time is the
timestamp of that triggering block. No transaction is needed to create
this derived event. An optional published audit record does not grant
it a second effect.

\begin{protocol}{alg:force}{Process the prelude and close a branch}
Before processing the payload of each authenticated finalized Espresso
block $J_q$, let $\theta=\theta_q$ and capture the entry epoch.
\begin{enumerate}[label=\textbf{\arabic*)},itemsep=3pt]
\item \textbf{Apply status.} From the complete control prefix, apply
all due independent STATUS and CLEAR records not yet applied, in
$(\theta_a,c)$ order. Mark them applied. Later blocks with the same
timestamp do not repeat them.

\item \textbf{Apply a multisignature Force.} Examine due Force
records in that same order. Apply only the first record whose target
matches the entry epoch. Capture the assigned export duty's context,
or the stored boundary context during Waiting. Apply that record's
attached configuration and close the epoch. Consume every other due
Force as stale, without applying its attached configuration. If a
matching Force closes the epoch, cancel its old export window and
skip step 3.

\item \textbf{Check delivery.} If an export session exists, settle
its deadlines at or below $\theta$ as in \cref{sec:windows}.
On the first failed window, capture that export duty's source region,
index and session ID,
set its cooldown to $\theta+C_{\rm auto}$, and close the epoch.
An unavailable namespace or control proof cannot establish a failure.

\item \textbf{Preserve and reopen.} For either closure, retain all
of $G$ and its provenance. Archive the old epoch's unconfirmed suffix
and normal successors. Record the cause, triggering header, retained
anchor, captured source context and selection lower bound $q$.
Create the new epoch. Do not accept its DATA in $J_q$ itself.

\item \textbf{Select or wait.} Record the pre-payload eligibility
snapshot $\chi_E(q)$. Resolve the new boundary, or an older Waiting
boundary, by its first eligible historical selection position.
When a duty opens, start a fresh startup budget at its canonical
opening timestamp. If no candidate exists, keep the anchor and run
no DATA window. Then return to payload processing.
\end{enumerate}
\end{protocol}

The cooldown is a replay-derived eligibility restriction. Governance ON
alone does not clear it; an explicit CLEAR can. Governance OFF continues
to exclude a region even after its cooldown expires.
Window and quota parameters remain fixed for the duty. A control that
needs to end it early uses an MSIG Force, not a retroactive deadline
change.

If Espresso time jumps across several scheduled thresholds, all those
events are first observable at the same block boundary. Step 1 processes
the independent status changes before either Force, and the matching
MSIG Force takes precedence over the entry session's automatic check.
No DATA is inserted between these prelude events. New startup windows
begin at the realized selection timestamp, not at an old missed
deadline, so a single block does not time out every new successor.

\subsection{Processing rules and control completeness}
\label{sec:control-completeness}
Control ingestion verifies a full L1 prefix and retains a pending
activation queue. A sequence number alone does not mean all ingested
records have activated. The interpreter processes each control once,
in its recorded target context. A closed epoch is never reopened by
an ON update, a repeated Force or a late regional receipt.

Every correct observer refreshes its complete finalized control view
before confirming an Espresso interval. The advance-announcement
assumptions ensure that all controls applicable to that interval are
already available. Missing evidence pauses replay. A new observer
reconstructs the same history; it does not reject an old instruction
merely because it learned it late.

A certifier's pure core must bind its supplied L1 checkpoint, control
sequence and history commitment. Its L1 acceptance interface must
reject an attestation that omits applicable controls. A conservative
implementation requires the advertised complete control version to
match the version required by the accepting contract; a stale version
is refreshed before retry. More selective checks require an equivalent
completeness proof. Proof of one event's inclusion is not proof that
no intervening control exists.

Discovery of a previously omitted control that should have affected
already confirmed history violates the completeness or announcement
premise. It is not a new permission to roll back $G$, apply the Force
at an arbitrary later point, or let nodes choose by local arrival time.
This protocol's conditional guarantees do not specify recovery from
that premise violation.

\section{The external observer}
\label{sec:observers}
An observer keeps the provisional production branch, persistent global
prefix, source-verification cursors and interpreter state. It does not
store per-block relay clocks or a committee-protected prefix.
The same deterministic logic can run inside the certifier.

\begin{protocol}{alg:external}{Verify Espresso output and advance the prefix}
\begin{enumerate}[label=\textbf{\arabic*)},itemsep=3pt]
\item \textbf{Authenticate inputs.} Obtain the next finalized
Espresso block and complete namespace evidence, linked to the trusted
quorum and prior prefix. Obtain the complete applicable finalized L1
control context and required historical authorities. If any required
history is unavailable, wait without treating it as empty.

\item \textbf{Run the prelude.} Execute \cref{alg:force} before
reading this block's payload. Persist the eligibility snapshot for this
Espresso position. Notify regional nodes of any derived forced boundary.

\item \textbf{Accept exports in order.} Decode records in canonical
namespace order. For each EXPORT, run \cref{sec:exports} on a temporary
state, including each normal boundary and session count. Commit the
whole update only if it passes. Invalid candidates do not change $G$
or restart a clock. Supporting evidence can be cached after validation.

\item \textbf{Persist and expose.} Persist the resulting state and
source cursors. Newly accepted DATA is global now, without another
timer or L1 transaction. Expose provisional output separately.
On closure, archive only the unconfirmed suffix and keep all of $G$.
\end{enumerate}
\end{protocol}

The expected base of an EXPORT is the DATA frontier and its binding
prefix digest, with provenance. It is not the hash of the entire
interpreter state: an empty Espresso block can change time, controls
and windows without changing this base. The certifier's state-transition
attestation, in contrast, binds the complete previous and next replay
checkpoints as well as the published range.

\subsection{Requests and retained data}
\label{sec:requests}
A request is \defterm{admissible} when it passes the application's admission
checks. Each node queues it in order of first receipt. Repeated delivery
does not change its queue position. New arrivals do not move
ahead of queued requests, and different nodes need not have the same
queue order. A correct leader skips requests
already in its proposed ancestry, including the inherited regional
prefix, and fills the batch with the oldest remaining requests that fit.
If no request can be included, it may propose an empty DATA block.

A request is finally removed from the pending queue only when included
in the global prefix. A request in a discarded regional suffix remains
pending and can be proposed on the new branch. Regional commitment or
regular handoff alone does not remove it. Keep the block data and
receipts needed for recovery and checkpoint verification; global
promotion does not authorize deleting their only available copy.



\paragraph{Faithful publication.}
The proposer or batcher publishes the next complete interval after
$C_{\rm pub}$, even if a Force has changed the current region.
The interval must belong to $G$, preserve order and multiplicity, and
match the DA payload and native execution evidence. An unavailable
accepted payload causes waiting, not an invented replacement.
The exact native adapter is specified separately: a Force cannot move
back from a partial-EVM anchor simply because the next signer cannot
continue it. It must obtain a legal continuation or wait.

\section{Examples and guarantees}
\label{sec:examples}
\subsection{Regular rotation and a late DRAIN}
\label{sec:regular-example}
Let $h(A_{r,0,0})=100$, $\TB_r=10$ and $K=4$. The normal DATA terms
produce 101--104, 105--108 and 109--110. A terminal QC below 110 moves
to the next term. A receipt for DATA 110 opens the next region, which
starts DATA 111 after Ready and bootstrap. No global wait is inserted.

\begin{figure}[htbp]\centering
% Protocols 3 and 4: a terminal receipt transfers the turn from A to B.
% Representative arrows only; time runs downward and is not to scale.
\begin{tikzpicture}[
  x=1cm,y=-1cm,font=\fontsize{9}{10.5}\selectfont,
  party/.style={draw,rounded corners=2pt,fill=white,align=center,
    inner sep=3pt,minimum width=2.2cm},
  lane/.style={draw=black!30,densely dashed},
  msg/.style={-{Stealth[length=1.7mm]},semithick,draw=rrBlue},
  backbone/.style={msg,densely dashed},
  tag/.style={fill=white,inner sep=1.5pt,font=\fontsize{8}{9}\selectfont},
  note/.style={anchor=north west,align=flush left,text width=5.55cm,
    inner sep=0pt,font=\fontsize{8.5}{10}\selectfont},
  phase/.style={anchor=west,fill=white,inner sep=2pt,
    font=\fontsize{9}{10.5}\selectfont\bfseries},
  dot/.style={circle,fill=black,inner sep=1.1pt},
  timer/.style={line width=2.6pt,draw=rrTeal!75},
  wait/.style={line width=2.6pt,draw=rrBlue!70!black,
    dash pattern=on 2.5pt off 1.5pt},
  entry/.style={draw=rrBlue,line width=1.2pt,dotted}
]
  \def\xA{1.2}\def\xR{3.8}\def\xL{6.4}\def\xB{9.0}\def\xn{10.0}
  \path[use as bounding box] (-.1,-.65) rectangle (15.6,15.75);
  \node[party] at (\xA,0) {A: leader /\\receipt holder};
  \node[party] at (\xR,0) {A: other\\replicas};
  \node[party] at (\xL,0) {B: initial leader\\$L(r+1,0,0)$};
  \node[party] at (\xB,0) {B: other\\replicas};
  \foreach \x in {\xA,\xR,\xL,\xB}
    \draw[lane] (\x,.55)--(\x,15.35);

  \node[phase] at (0,.85) {1.\ A commits its quota endpoint};
  \draw[msg] (\xA,1.40)--node[tag,above,sloped] {$T_r=110$}(\xR,1.70);
  \draw[msg] (\xR,2.15)--node[tag,above,sloped] {DATA votes}(\xA,2.45);
  \foreach \x/\s/\e in {\xA/1.40/3.80,\xR/1.70/4.10} {
    \draw[timer] (\x+.20,\s)--(\x+.20,\e);
    \node[dot] at (\x,\s) {};
    \draw[thick] (\x-.12,\e)--(\x+.45,\e);
  }
  \node[note] at (\xn,1.35) {Replicas validate and vote for the last DATA
    block. Each starts its own $2\Delta_{\rm A}$ commit timer.};
  \node[tag,anchor=east] at (\xA-.15,2.90) {QC$(110)$};
  \node[note] at (\xn,2.80) {The terminal QC does not hand the turn to B.
    The quota permits no DATA 111 in A's session.};
  \draw[msg] (\xR,4.25)--node[tag,above,sloped] {FINAL$(110)$}(\xA,4.55);
  \node[dot] at (\xA,4.55) {};
  \node[note] at (\xn,3.95) {Timers expire; replicas commit 110 and send
    FINAL. Matching signatures of weight $>F_{\rm A}$ form
    $\FC_r(T_r)$. Receipt holders complete the old engine.};

  \node[phase] at (0,5.55) {2.\ The regular boundary crosses the backbone};
  \node[note] at (\xn,5.50) {Verify the authorized session and quota.
    Derive B from $\chi_E(q_s)$ and form the REGULAR witness;
    no new L1 transaction is needed.};
  \draw[backbone] (\xA,6.35)--node[tag,pos=.55,above,sloped]
    {REGULAR witness + bundle}(\xL,7.35);
  \draw[backbone] (\xA,6.35)--(\xB,7.55);
  \node[note] at (\xn,6.90) {Transfer has no fixed delay bound. Global
    confirmation of 110 proceeds separately; this handoff does not wait
    for it.};

  \node[phase] at (0,8.05) {3.\ B validates, certifies the anchor, and produces};
  \draw[entry] (\xL-.6,8.45)--(\xL,8.45)--(\xB,8.65)--(\xB+.6,8.65);
  \foreach \x/\y in {\xL/8.45,\xB/8.65} \node[dot] at (\x,\y) {};
  \draw[msg] (\xL,8.75)--node[tag,above,sloped] {Ready relay}(\xB,8.95);
  \draw[msg] (\xB,9.25)--node[tag,above,sloped] {STATUS}(\xL,9.50);
  \draw[wait] (\xL+.20,8.45)--(\xL+.20,10.85);
  \node[note] at (\xn,8.40) {Ready requires the boundary, anchor 110,
    history, and finalized authority. Enter $(r+1,0,0)$; relay the bundle
    and send STATUS.};
  \node[note] at (\xn,9.95) {The initial leader waits $2\Delta_{\rm B}$
    from its own entry. Progress checks already run.};
  \draw[msg] (\xL,10.90)--node[tag,above,sloped] {NEWVIEW$(110)$}(\xB,11.20);
  \draw[msg] (\xB,11.65)--node[tag,above,sloped] {ANCHOR votes}(\xL,11.95);
  \node[dot] at (\xL,11.95) {};
  \node[note] at (\xn,11.05) {NEWVIEW must choose the inherited anchor.
    Weight $>F_{\rm B}$ gives an ANCHOR QC; these votes start no timer.};
  \draw[msg] (\xL,12.55)--node[tag,above,sloped] {DATA 111}(\xB,12.85);
  \foreach \x/\s/\e in {\xL/12.55/14.95,\xB/12.85/15.25} {
    \draw[timer] (\x+.20,\s)--(\x+.20,\e);
    \node[dot] at (\x,\s) {};
    \draw[thick] (\x-.12,\e)--(\x+.45,\e);
  }
  \draw[msg] (\xB,13.35)--node[tag,above,sloped] {DATA votes}(\xL,13.65);
  \node[note] at (\xn,12.55) {The first fresh block has DATA height 111
    and physical position 1. Its votes start new commit timers. A QC
    permits the next proposal before 111 commits.};
  \node[note] at (\xn,14.60) {If uninterrupted, the new timers commit 111;
    replicas send FINAL for B's session.};
\end{tikzpicture}

\caption{Normal cross-region rotation retains the DATA, FINAL and
ANCHOR sequence. The historical Espresso selection snapshot fixes the
successor. Regional production and Espresso publication proceed in
parallel. Neither the terminal receipt nor the new region waits for
global confirmation.}
\label{fig:regular-cross-region}
\end{figure}
\FloatBarrier

Now suppose 110 has a QC but its timer is cancelled by Quit.
Recovery opens at height 110 and proposes a DRAIN, which supplies a
new timer without using DATA quota. When it commits the terminal DATA,
replicas send FINAL for 110, not for DRAIN. The normal successor still
inherits 110. Its old L1 and Espresso references do not establish an
earlier commitment or export time.

Suppose $G$ is only at 100 while A has committed 110 and B has committed
111. A publisher can export 101--111 with the required A and B receipts
and normal witness. Once that complete EXPORT is accepted from finalized
Espresso history, the entire interval becomes global. A lone record for
111 cannot bypass missing ancestors.

\subsection{A multisignature Force and the retained anchor}
\label{sec:force-example}
Suppose the committee submits a Force for epoch $\beta$ at real time
$s$. It selects $\theta_a$ with the stated announcement margin. All
correct participants can verify the finalized record before the first
Espresso block whose timestamp reaches $\theta_a$ becomes final.
Before that block, valid old-epoch exports can still advance $G$.

If $G$ advances from 100 to 108 during this interval, the Force retains
108. It cannot retain the older observation 100. At the crossing block,
the interpreter closes the old epoch before processing payloads.
Later old-epoch DATA 109 is not accepted, even if it has a valid regional
receipt. The selected successor uses a new identity and continues at
109 after Ready. No Espresso-to-L1 checkpoint update is needed to start
this handoff.

\begin{figure}[htbp]\centering
% Two Force causes share one Espresso boundary and the existing regional bootstrap.
% Vertical spacing illustrates event order, not elapsed time.
\begin{tikzpicture}[
  x=1cm,y=-1cm,font=\fontsize{8.3}{9.8}\selectfont,
  party/.style={draw,rounded corners=2pt,fill=white,align=center,
    inner sep=3pt,text width=1.22cm,minimum height=.66cm},
  lane/.style={draw=black!28,densely dashed},
  msg/.style={-{Stealth[length=1.6mm]},semithick,draw=rrBlue},
  read/.style={msg,densely dashed},
  force/.style={msg,draw=rrRed},
  tag/.style={fill=white,inner sep=1.5pt,font=\fontsize{7.7}{9}\selectfont},
  note/.style={anchor=north west,align=flush left,text width=5.55cm,
    inner sep=0pt,font=\fontsize{8.3}{9.8}\selectfont},
  phase/.style={anchor=west,fill=white,inner sep=2pt,
    font=\fontsize{9}{10.5}\selectfont\bfseries},
  dot/.style={circle,fill=black,inner sep=1.1pt},
  wait/.style={line width=2.5pt,draw=rrTeal!75,
    dash pattern=on 2.5pt off 1.5pt},
  rulebox/.style={draw=rrRed!55,rounded corners=2pt,fill=rrRed!3,
    align=center,text width=7.9cm,inner sep=5pt}
]
  \def\xC{.75}\def\xI{2.3}\def\xE{3.85}
  \def\xN{5.4}\def\xL{6.95}\def\xB{8.5}\def\xn{9.45}
  \path[use as bounding box] (-.05,-.55) rectangle (15.05,16.05);
  \node[party] at (\xC,0) {Multisig};
  \node[party] at (\xI,0) {L1};
  \node[party] at (\xE,0) {Espresso};
  \node[party] at (\xN,0) {Node /\\certifier};
  \node[party] at (\xL,0) {B: initial\\leader};
  \node[party] at (\xB,0) {B: other\\replicas};
  \foreach \x in {\xC,\xI,\xE,\xN,\xL,\xB}
    \draw[lane] (\x,.52)--(\x,15.65);

  \node[phase] at (0,.88) {1.\ Announce MSIG Force on L1};
  \draw[force] (\xC,1.55)--(\xI,1.85);
  \node[tag,anchor=south] at (1.52,1.48) {submit $F$ at $s$};
  \node[note] at (\xn,1.30) {$F$ names old branch epoch $\beta$ and
    activation time $\theta_a$. It does not choose an older retained
    DATA height.};
  \node[dot] at (\xI,2.65) {};
  \draw[read,draw=rrRed] (\xI,2.65)--node[tag,above,sloped]
    {finalized control + complete history}(\xN,3.30);
  \node[dot] at (\xN,3.30) {};
  \node[note] at (\xn,2.85) {By $s+T_{\rm L1}$, all correct participants
    can verify the finalized control. $T_{\rm L1}$ includes dissemination
    and verification. This precedes the Espresso crossing.};
  \node[tag,align=center] at (5.35,4.15)
    {Before activation: old-epoch output remains eligible;\\
     normal regional handoffs may continue.};
  \node[phase] at (0,5.25) {2.\ Finalize the Espresso boundary; derive one Force};
  \draw[force] (\xE,5.95)--(\xN,6.25);
  \node[tag,anchor=south] at (4.62,5.90) {$J_{q^*}$ finalized};
  \node[note] at (\xn,5.75) {For MSIG Force,
    $q^*=\min\{q:\theta_q\ge\theta_a\}$.
    Interpret the complete history through $J_{q^*-1}$ first.};
  \node[rulebox] at (4.65,7.30) {At the start of $J_{q^*}$:\\[2pt]
    due independent status changes $\;\to\;$ first matching MSIG Force;\\
    otherwise check due windows and derive AUTO Force;\\
    only then process EXPORT records.};
  \node[note] at (\xn,7.05) {AUTO is the alternative cause: fewer than
    $m_{\rm start}$ or $q_{\min}$ DATA were accepted in a fixed window.
    The crossing triggers closure without an L1 transaction.};
  \node[dot,fill=rrRed] at (\xN,8.55) {};
  \node[tag,anchor=east,text=rrRed] at (\xN-.12,8.55)
    {anchor $:=G$ before payload};
  \node[note] at (\xn,8.70) {Close old $\beta$ and its speculative
    regular successors. Keep all of $G$. Select eligible B;
    if no successor exists, wait.};
  \draw[read] (\xN,9.60)--node[tag,pos=.55,above,sloped]
    {complete boundary}(\xL,10.00);
  \draw[read] (\xN,9.60)--(\xB,10.15);
  \node[note] at (\xn,10.15) {B verifies both chains, the anchor,
    authority and state before Ready. New-epoch output is accepted
    only after the triggering block.};

  \node[phase] at (0,11.10) {3.\ B uses the existing ANCHOR bootstrap};
  \node[dot] at (\xL,11.65) {};
  \node[dot] at (\xB,11.85) {};
  \draw[msg] (\xB,12.20)--node[tag,above,sloped] {STATUS}(\xL,12.45);
  \draw[wait] (\xL+.18,11.65)--(\xL+.18,13.15);
  \node[tag,anchor=east] at (\xL-.12,12.85) {wait $2\Delta_{\rm B}$};
  \node[note] at (\xn,11.75) {Install a fresh session ID and the exact
    anchor. The leader waits $2\Delta_{\rm B}$ after its own Ready,
    then sends NEWVIEW.};
  \draw[msg] (\xL,13.20)--node[tag,above,sloped] {NEWVIEW}(\xB,13.45);
  \draw[msg] (\xB,14.00)--node[tag,above,sloped] {ANCHOR votes}(\xL,14.25);
  \node[dot] at (\xL,14.25) {};
  \node[note] at (\xn,13.65) {An ANCHOR QC with weight $>F_{\rm B}$
    permits DATA after $G$. ANCHOR votes start no commit timer.
    DATA and DRAIN retain their rules.};
  \draw[msg] (\xL,15.00)--node[tag,above,sloped]
    {DATA at $h(G)+1$}(\xB,15.30);
  \node[note] at (\xn,15.20) {The forced handoff preserves $G$;
    it does not wait for $G$ to be registered on L1.};
\end{tikzpicture}

\caption{Multisignature Force: L1 supplies a future-dated instruction;
the first finalized Espresso timestamp crossing fixes its boundary.
The complete pre-boundary global prefix survives. State transfer and
the original regional bootstrap follow; old regional engines are not
silently stopped.}
\label{fig:forced-cross-region}
\end{figure}
\FloatBarrier

\subsection{A failed window and equal timestamps}
\label{sec:window-example}
For illustration, let the startup window be $[0,4)$ with target one,
and later windows have length four and target two, subject to remaining
quota. One DATA at time 1 satisfies startup. If only one further DATA
is accepted at time 5, the running window $[4,8)$ fails.
A record in the first block timestamped 8 is too late to repair it.
The automatic Force preserves the two accepted DATA and closes their
epoch before reading that block's payload.

\begin{figure}[htbp]\centering
% Shared replay semantics: normal production runs ahead; a deadline is exclusive.
% Panel (b) is a separate numerical example, with Espresso time in fixed units.
\begin{tikzpicture}[
  x=1cm,y=-1cm,font=\fontsize{8.7}{10.2}\selectfont,
  box/.style={draw=rrBlue!65,rounded corners=2pt,fill=white,
    align=center,inner sep=5pt,text width=3.80cm,minimum height=.90cm},
  outcome/.style={box,draw=rrTeal,fill=rrTeal!4},
  warning/.style={box,draw=rrRed,fill=rrRed!3},
  msg/.style={-{Stealth[length=1.7mm]},semithick,draw=rrBlue},
  tag/.style={fill=white,inner sep=2pt,font=\fontsize{8}{9.3}\selectfont},
  phase/.style={anchor=west,fill=white,inner sep=2pt,
    font=\fontsize{9.2}{10.7}\selectfont\bfseries},
  note/.style={anchor=north west,align=flush left,inner sep=0pt,
    font=\fontsize{8.7}{10.2}\selectfont},
  dot/.style={circle,fill=rrTeal,inner sep=2pt},
  cross/.style={draw=rrRed,line width=1.1pt}
]
  \path[use as bounding box] (0,0) rectangle (15.15,13.85);
  \node[phase] at (.05,.30) {(a)\ Normal handoff and global confirmation use separate paths};
  \node[box] (terminal) at (2.15,2.00)
    {A's terminal DATA $T_A$\\FINAL + full state};
  \node[outcome] (ready) at (7.55,1.35)
    {REGULAR witness verified\\B Ready\\ANCHOR bootstrap};
  \node[outcome] (produce) at (12.95,1.35)
    {B starts new DATA\\No global-finality wait};
  \node[box] (publish) at (7.55,3.12)
    {Continuous EXPORT\\Each session's evidence};
  \node[outcome] (global) at (12.95,3.12)
    {Espresso finalized\\Replay accepts DATA into $G$};
  \draw[msg] (terminal.east)--(ready.west);
  \draw[msg] (ready.east)--(produce.west);
  \draw[msg] (terminal.east)--(publish.west);
  \draw[msg] (publish.east)--(global.west);
  \node[note,text width=14.6cm] at (.25,4.15)
    {Only DATA accepted into $G$ counts toward the export window. Regional
     production may already be in later sessions; DRAIN, QC, duplicates
     and raw inclusion do not count.};

  \node[phase] at (.05,5.20) {(b)\ AUTO example: one accepted DATA in $[4,8)$, with $q_{\min}=2$};
  \draw[line width=4pt,draw=rrTeal!22] (1.25,6.80)--(10.45,6.80);
  \draw[draw=rrTeal,line width=1pt] (1.25,6.55)--(1.25,7.05);
  \draw[draw=rrTeal,fill=white,line width=1pt] (10.45,6.80) circle (.075);
  \draw[msg,draw=black!65] (.60,7.20)--(14.65,7.20)
    node[tag,above,anchor=east] {Espresso timestamp};
  \foreach \x/\t in {1.25/4,3.55/5,5.85/6,8.15/7,10.45/8,12.75/9} {
    \draw[black!65] (\x,7.10)--(\x,7.30);
    \node[tag,anchor=north] at (\x,7.34) {$\t$};
  }
  \node[dot] at (3.55,6.80) {};
  \node[tag,align=center,anchor=south] at (3.55,6.51)
    {one accepted DATA\\at $\theta=5$};
  \draw[draw=rrRed,densely dashed,semithick] (10.45,5.90)--(10.45,9.55);
  \node[tag,anchor=west,text=rrRed] at (10.65,6.28)
    {first block with $\theta_q\ge8$};
  \node[warning,text width=7.5cm] (failure) at (4.90,8.75)
    {At block start: accepted count $1<2$.\\
     No matching MSIG Force: apply AUTO.\\Keep all of $G$.};
  \draw[cross] (10.33,8.35)--(10.57,8.59);
  \draw[cross] (10.33,8.59)--(10.57,8.35);
  \node[note,text width=4.00cm,text=rrRed] at (10.78,8.18)
    {Old-epoch DATA at $\theta=8$ is too late. Its raw inclusion cannot
     repair the expired window or append to $G$.};
  \node[note,text width=14.6cm] at (.25,9.75)
    {The interval is half-open. All blocks with timestamp $8$ lie outside
     $[4,8)$. A later block with the same timestamp cannot undo the first
     crossing. AUTO needs no L1 call or separately published Force tag.};

  \node[phase] at (.05,10.85) {(c)\ Every observer and certifier uses the same block-start order};
  \node[box,text width=3.0cm,minimum height=.96cm] (status) at (1.85,11.85)
    {Due independent\\status updates};
  \node[warning,text width=3.0cm,minimum height=.96cm] (msig) at (5.62,11.85)
    {First matching\\MSIG Force};
  \node[warning,text width=3.0cm,minimum height=.96cm] (auto) at (9.39,11.85)
    {Otherwise: due windows\\and AUTO if failed};
  \node[outcome,text width=3.0cm,minimum height=.96cm] (data) at (13.16,11.85)
    {Then validate\\EXPORT records};
  \draw[msg] (status.east)--(msig.west);
  \draw[msg] (msig.east)--(auto.west);
  \draw[msg] (auto.east)--(data.west);
  \node[note,text width=14.6cm] at (.25,12.88)
    {Both Forces close future appends by the old branch epoch and preserve
     all of $G$. A selected forced successor receives a fresh startup
     window. Missing authenticated history makes the observer wait; it is
     not evidence of an empty namespace.};
\end{tikzpicture}

\caption{External confirmation and a block-missing Force. Only accepted
DATA counts. The running-window example assumes at least two DATA
remain due; in general the target is capped by remaining quota.
Deadline-equal output cannot revive the expired epoch.}
\label{fig:observer-finality}
\end{figure}
\FloatBarrier

If Espresso jumps from timestamp 5 to 30, the same first failed window
is detected at the latter block. A new successor starts its window at
30, not at 8. If another Espresso block also has timestamp 30, it does
not repeat the old control events or Force. A due status update at the
same crossing is applied first, so the successor uses the prescribed
eligibility snapshot.

\subsection{Conditional prefix-safety argument}
\label{sec:safety-argument}
Assume the source-history, announcement, regional validity and
verification premises above. Start two correct observers from the same
trusted interpreter checkpoint. We establish the following properties
by induction over complete Espresso blocks and their ordered records.

\paragraph{Common interpretation.}
The observers authenticate compatible Espresso prefixes. Every control
applicable at a block's timestamp is already available by the
announcement condition, and each observer checks the complete control
context. Due-record ordering, historical selection, fixed windows and
payload order are deterministic. If an observer lacks evidence it waits;
it does not choose a different candidate or infer an empty interval.
Thus both obtain the same state whenever they have verified the same
prefix. A slower observer's global DATA is a prefix of the faster one's.

\paragraph{Preservation of confirmed DATA.}
An accepted EXPORT adds exactly the next continuous interval. Rejected
records add nothing. Independent status updates do not change DATA.
Either Force takes its anchor from the current $G$ and changes only the
authority for subsequent output. Waiting and window transitions also
remove no DATA. Therefore each transition preserves the entire preceding
global prefix and its provenance.

\paragraph{No re-entry of the closed suffix.}
All outputs bind a session and its branch epoch. After closure, the
acceptance rule rejects new old-epoch DATA, including normal successors
that started before the Force. A new session may reuse a numeric index
or DATA height beyond $G$, but its identity is different. It cannot
replace any retained block. Old regional receipts remain true statements
about old commitments without authorizing a new global append.

The argument does not re-prove the regional DATA/DRAIN engine or assert
that every regional result is eventually delivered. It also does not
make malicious or late control announcement harmless: non-retroactive
announcement and complete verification are stated premises. Sound
execution and DA publication remain separate conditions.

\subsection{Progress conditions and limits}
The regional engine keeps its original fair-schedule, exclusion and
uninterrupted-commit-window conditions. Normal handoff still needs
state transfer and authority, but no Espresso confirmation of its
terminal DATA. Global output additionally needs publisher service,
Espresso inclusion, and enough sustained delivery to satisfy the fixed
windows. Eventual delivery alone does not give this rate guarantee.

A Force needs a recoverable retained anchor and an eligible successor.
Its startup budget starts at the public opening, not at a replica's
later Ready time. Fixed budgets therefore also assume timely startup
for some serving interval. An empty DATA can satisfy protocol delivery
without establishing user-transaction inclusion or censorship resistance.
With only one eligible region, the unchanged rule that selects another
region can wait indefinitely. Same-region renewal is not introduced here.

Espresso's own failure prevents new global confirmation and automatic
crossings. A future-dated multisignature Force removes the need for a
second L1 round trip, but still waits for its announced Espresso
boundary and subsequent Ready. This document gives no unconditional
wall-clock bound on failover or withdrawal.

\subsection{Open design and implementation questions}
\begin{enumerate}[label=(\arabic*)]
\item Pin the concrete Espresso finality, timestamp and namespace-proof
profile, including clock lead and history retention.
\item Fix registry/control-proof encodings and the certifier's
complete-control-version acceptance rule.
\item Choose window and cooldown parameters from explicit service
assumptions, and decide whether a stronger normal-reference freshness
rule is needed.
\item Specify native execution and signing continuation at every
allowed global anchor, including partial-EVM state if supported.
\item Decide when a closed old engine may stop before Completion,
without treating the external closure as a regional quit proof.
\item Complete staking retirement and publisher incentives. A missed
delivery window alone is not a proof of proposer equivocation.
\end{enumerate}

\begingroup\small
\begin{thebibliography}{9}
\bibitem{part1}
Sync HotStuff RR, Part I: Core Protocol. Design document, 13 October 2026.
Companion source: \href{run:sync_hotstuff_rr_part1_core.tex}{Part I core protocol}.
\bibitem{sync}
I. Abraham, D. Malkhi, K. Nayak, L. Ren, and M. Yin.
\emph{Sync HotStuff: Simple and Practical Synchronous State Machine Replication}.
IEEE Symposium on Security and Privacy, 2020, pp. 106--118.
\href{https://conferences.computer.org/sp/pdfs/sp/2020/349700a106.pdf}{Published paper}.
\bibitem{certification}
MegaETH Specification, PR~8, \emph{Block Certification}.
Commit \texttt{aedc2826663b}.
\href{https://github.com/megaeth-labs/specification/tree/aedc2826663b134ee25d26bea17e70c6d23d321c/specs/block-certification}{Pinned specification}.
Its Espresso and certifier verification interfaces are a reference;
the rotation and window rules in this document are new.
\end{thebibliography}
\endgroup
\end{document}
